% !TEX root = TI5.tex

\section{Faster algorithms for isomorphism of odd order class-2 nilpotent groups}
 
%\begin{theorem}
%[Main result for \algprobm{Group Isomorphism}]
%There is an algorithm to de
%\end{theorem}

In this section, we prove Theorem~\ref{thm:tensor_isomorphism_main} by giving an algorithm for testing isomorphism of two nilpotent class-2 groups of odd order in time $N^{O((\log N)^{8/9})}$, where $N$ is the order of the two groups.

In Section~\ref{sec:group_iso_to_tensor_iso}, we reduce \algprobm{Group Isomorphism} for class-2 nilpotent groups to \algprobm{Alternating Matrix Space Isometry} over $\Z/p^e \Z$ for some prime $p$ and integer $e$. 
In Section~\ref{sec:tensor_iso_to_matrix_tuple}, we further reduce \algprobm{Alternating Matrix Space Isometry} to \algprobm{Alternating Matrix Tuple Isometry}. 
In Section~\ref{sec:matrix_tuple}, we give a polynomial time algorithm for the \algprobm{Alternating Matrix Tuple Isometry} problem. 

In this section, we say that
an $n \times n \times m$ tensor $\tA = (a_{i, j, k})$ is alternating (between the first and second directions) if, for any $1 \leq i, j \leq n$ and $k \leq m$, $a_{i, j, k} = -a_{j, i, k}$. 
The \algprobm{Alternating Matrix Space Isometry} problem %, the input contains two $n \times n \times m$ alternating tensors $\tA$ and $\tB$ over $\Z / p^e \Z$ for some prime $p$ and positive integers $n, m, e$, 
is to determine  whether for two $n \times n \times m$ alternating tensors $\tA$ and $\tB$ over $\Z / p^e \Z$, there are two matrices $P \in \GL(n, \Z / p^e \Z)$ and $Q \in \GL(m, \Z / p^e \Z)$ such that $\tB = (P, P, Q) \circ \tA$.
The \algprobm{Alternating Matrix Tuple Isometry} problem %, the input contains two $n \times n \times m$ alternating tensors $\tA$ and $\tB$ over $\Z / p^e \Z$ for some prime $p$ and positive integers $n, m, e$, 
is to determine whether for two alternating matrix tuples $(A_1, \dots, A_k)$ and $(B_1, \dots, B_k)$ in $\M(n, \Z / p^e \Z)$, there is a matrix $P \in \GL(n, \Z / p^e \Z)$ such that $B_i = P A_i P^t$ for any $1 \leq i \leq k$.


\subsection{Reduction to \algprobm{Alternating Matrix Space Isometry} over $\Z/p^e \Z$}\label{sec:group_iso_to_tensor_iso}
%\ynote{This subsection needs to be filled in. I assume that the reduction would be to isomorphism problems for $\Lambda^2((\Z/p^e\Z)^n)\otimes (\Z/p^f\Z)^m$. Is this understanding correct?} \jnote{Now it is :).}

%\begin{theorem}[Baer's Correspondence]
%
%\end{theorem}

%An abelian group is called \emph{homocyclic} if it is of the form $(\Z/m \Z)^d$ for some $m,d$. By extension, we call a Lie ring homocyclic if its underlying Abelian group is homocyclic.
%
%\begin{lemma} \label{lem:homocyclic}
%There is a reduction from \algprobm{Lie Ring Isomorphism} for Lie rings over $\Z/p^e \Z$ to \algprobm{Homocyclic Lie Ring Isomorphism} over $\Z/p^e \Z$.
%\end{lemma}
%
%\begin{proof}
%As $\Z/p^e \Z$-modules, $\Z/p^k \Z$ is isomorphic to $p^{e-k} \Z / p^e \Z$. Thus, if the underlying Abelian group of a Lie ring is $\bigoplus_{i=1}^n \Z/p^{e_i} \Z$, then 
%\end{proof}
\begin{comment}
\begin{definition}\label{def:tensor_from_group}
Let $G$ and $H$ be two Abelian $p$-groups for some odd prime $p$, and $\tA$ be a tensor in $G \otimes G\otimes H$ with $f_\tA$ as the multilinear function associated with $\tA$. 
Let $e$ be an integer such that the largest order of elements in both $G$ and $H$ is at most $p^e$. 
Suppose $G$ is isomorphic to 
\[\bigoplus_{\ell = 1}^{e} \left(\Z / p^{\ell } \Z\right)^{\alpha_\ell},\]
and $H$ is isomorphic to 
\[\bigoplus_{\ell = 1}^{e} \left(\Z / p^{\ell } \Z\right)^{\beta_\ell}.\]
Let $n = \sum_{\ell = 1}^{e} \alpha_\ell$ and $m = \sum_{\ell = 1}^{e} \beta_\ell $. 
An $n \times n \times m$ tensor $\tB = (b_{i, j, k})$ over $\Z / p^e \Z$ is a tensor derived from $\tA$ if there exist
$c_1, \dots, c_n \in G$ and $d_1, \dots, d_m \in H$ satisfying the following conditions:
\begin{enumerate}
\item $c_1, \dots, c_n$ generate $G$ and for each $1 \leq \ell \leq e$, $c_{j}$ is of order $p^\ell$ for any $(\sum_{i = 1}^{\ell - 1} \alpha_i) + 1 \leq j \leq \sum_{i = 1}^{\ell} \alpha_i$. 
\item $d_1, \dots, d_m$ generate $H$ and for each $1 \leq \ell \leq e$, $d_{j}$ is of order $p^\ell$ for any $(\sum_{i = 1}^{\ell - 1} \beta_i) + 1 \leq j \leq \sum_{i = 1}^{\ell} \beta_i$. 
\item For any $1 \leq i, j \leq n, 1 \leq k \leq m$, $b_{i, j, k} = f_\tA(c_i, c_j, d_k)\cdot p^{e - t}$, where $t$ is the integer such that the order of $d_k$ is $p^t$. 
\end{enumerate}



%Suppose $[G, G]$ is isomorphic to 
%\[\bigoplus_{\ell = 1}^{e} \left(\Z / p^{\ell } \Z\right)^{\alpha_\ell},\]
%and $G / Z(G)$ is isomorphic to 
%\[\bigoplus_{\ell = 1}^{e} \left(\Z / p^{\ell } \Z\right)^{\beta_\ell}.\]
%Let $m = \sum_{\ell = 1}^{e} \alpha_\ell$ and $n = \sum_{\ell = 1}^{e} \beta_\ell $. 



%and $e$ be the smallest integer such that the largest order of elements in both $[G, G]$ and $G / Z(G)$ is at most $p^e$. 
%Suppose $[G, G]$ is isomorphic to 
%\[\bigoplus_{\ell = 1}^{e} \left(\Z / p^{\ell } \Z\right)^{\alpha_\ell},\]
%and $G / Z(G)$ is isomorphic to 
%\[\bigoplus_{\ell = 1}^{e} \left(\Z / p^{\ell } \Z\right)^{\beta_\ell}.\]
%Let $m = \sum_{\ell = 1}^{e} \alpha_\ell$ and $n = \sum_{\ell = 1}^{e} \beta_\ell $. 

%An $n \times n \times m$ tensor $\tA = (a_{i, j, k})$ over $\Z / p^e \Z$ is a tensor derived from $G$ if there exist
%$c_1, \dots, c_n \in G / Z(G)$ and $d_1, \dots, d_m$ satisfying the following conditions:
%\begin{enumerate}
%\item $c_1, \dots, c_n$ generate $G / Z(G)$ and for each $1 \leq \ell \leq e$, $c_{j}$ is of order $p^\ell$ for any $(\sum_{i = 1}^{\ell - 1} \beta_i) + 1 \leq j \leq \sum_{i = 1}^{\ell} \beta_i$. 
%\item $d_1, \dots, d_m$ generate $[G, G]$ and for each $1 \leq \ell \leq e$, $d_{j}$ is of order $p^\ell$ for any $(\sum_{i = 1}^{\ell - 1} \alpha_i) + 1 \leq j \leq \sum_{i = 1}^{\ell} \alpha_i$. 
%\item For any $1 \leq i, j \leq n$, $[c_i, c_j] = \prod_{k = 1}^md_k^{a_{i, j, k}} \cdot p^{e - t}$, where $t$ is the integer such that the order of $d_k$ is $p^t$. 
%\end{enumerate}
\end{definition}



\begin{definition}\label{def:matrix_from_group}
Let $G$ and $H$ be two Abelian $p$-groups for some odd prime $p$, and $A$ be a matrix in $G \otimes H$ with $f_A$ as the bilinear function associated with $A$. 
Let $e$ be an integer such that the largest order of elements in both $G$ and $H$ is at most $p^e$. 
Suppose $G$ is isomorphic to 
\[\bigoplus_{\ell = 1}^{e} \left(\Z / p^{\ell } \Z\right)^{\alpha_\ell},\]
and $H$ is isomorphic to 
\[\bigoplus_{\ell = 1}^{e} \left(\Z / p^{\ell } \Z\right)^{\beta_\ell}.\]
Let $n = \sum_{\ell = 1}^{e} \alpha_\ell$ and $m = \sum_{\ell = 1}^{e} \beta_\ell $. 
An $n \times m$ matrix  $B = (b_{i, j})$ over $\Z / p^e \Z$ is a matrix derived from $A$ if there exist
$c_1, \dots, c_n \in G$ and $d_1, \dots, d_m \in H$ satisfying the following conditions:
\begin{enumerate}
\item $c_1, \dots, c_n$ generate $G$ and for each $1 \leq \ell \leq e$, $c_{j}$ is of order $p^\ell$ for any $(\sum_{i = 1}^{\ell - 1} \alpha_i) + 1 \leq j \leq \sum_{i = 1}^{\ell} \alpha_i$. 
\item $d_1, \dots, d_m$ generate $H$ and for each $1 \leq \ell \leq e$, $d_{j}$ is of order $p^\ell$ for any $(\sum_{i = 1}^{\ell - 1} \beta_i) + 1 \leq j \leq \sum_{i = 1}^{\ell} \beta_i$. 
\item For any $1 \leq i, j \leq n, 1 \leq k \leq m$, $b_{i, j} = f_A(c_i, d_j)\cdot p^{e - t}$, where $t$ is the integer such that the order of $d_j$ is $p^t$. 
\end{enumerate}



%Suppose $[G, G]$ is isomorphic to 
%\[\bigoplus_{\ell = 1}^{e} \left(\Z / p^{\ell } \Z\right)^{\alpha_\ell},\]
%and $G / Z(G)$ is isomorphic to 
%\[\bigoplus_{\ell = 1}^{e} \left(\Z / p^{\ell } \Z\right)^{\beta_\ell}.\]
%Let $m = \sum_{\ell = 1}^{e} \alpha_\ell$ and $n = \sum_{\ell = 1}^{e} \beta_\ell $. 



%and $e$ be the smallest integer such that the largest order of elements in both $[G, G]$ and $G / Z(G)$ is at most $p^e$. 
%Suppose $[G, G]$ is isomorphic to 
%\[\bigoplus_{\ell = 1}^{e} \left(\Z / p^{\ell } \Z\right)^{\alpha_\ell},\]
%and $G / Z(G)$ is isomorphic to 
%\[\bigoplus_{\ell = 1}^{e} \left(\Z / p^{\ell } \Z\right)^{\beta_\ell}.\]
%Let $m = \sum_{\ell = 1}^{e} \alpha_\ell$ and $n = \sum_{\ell = 1}^{e} \beta_\ell $. 

%An $n \times n \times m$ tensor $\tA = (a_{i, j, k})$ over $\Z / p^e \Z$ is a tensor derived from $G$ if there exist
%$c_1, \dots, c_n \in G / Z(G)$ and $d_1, \dots, d_m$ satisfying the following conditions:
%\begin{enumerate}
%\item $c_1, \dots, c_n$ generate $G / Z(G)$ and for each $1 \leq \ell \leq e$, $c_{j}$ is of order $p^\ell$ for any $(\sum_{i = 1}^{\ell - 1} \beta_i) + 1 \leq j \leq \sum_{i = 1}^{\ell} \beta_i$. 
%\item $d_1, \dots, d_m$ generate $[G, G]$ and for each $1 \leq \ell \leq e$, $d_{j}$ is of order $p^\ell$ for any $(\sum_{i = 1}^{\ell - 1} \alpha_i) + 1 \leq j \leq \sum_{i = 1}^{\ell} \alpha_i$. 
%\item For any $1 \leq i, j \leq n$, $[c_i, c_j] = \prod_{k = 1}^md_k^{a_{i, j, k}} \cdot p^{e - t}$, where $t$ is the integer such that the order of $d_k$ is $p^t$. 
%\end{enumerate}
\end{definition}


\begin{lemma}
Let $G, H, G', H'$ be four Abelian $p$-groups such that $G \cong G'$ and $H \cong H'$, and $\tA, \tA'$ be two tensors in $G \otimes G  \otimes H$ and $G' \otimes G' \otimes H'$ respectively.  
Let $\tB$ and $\tB'$ be two tensors over $\Z / p^e \Z$ for some integer $e$ derived from $\tA$ and $\tA'$ respectively. 
Then $\tA$ and $\tA'$ are isomorphic if and only if their derived tensors are isomorphic.
\end{lemma}
\begin{proof}
By Definition~\ref{def:tensor_from_group}, the two tensors are of the same dimension. 
Let $c_{\tA, 1}, \dots, c_{\tA, n}, d_{\tA, 1}, \dots, d_{\tA, m}$ and $c_{\tA', 1}, \dots, c_{\tA', n}, d_{\tA', 1}, \dots, d_{\tA', m}$ be the elements used to obtain $\tB$ and $\tB'$  following Definition~\ref{def:tensor_from_group} respectively. 

Since both $\{c_{\tA, 1}, \dots, c_{\tA, n}\}$ and $\{c_{\tB, 1}, \dots, c_{\tB, n}\}$ are generating sets of $G$, there are $x_{i, j}$ such that $c_{\tB, i} = \prod_{j = 1}^n c_{\tA, j}^{x_{i, j}}$ for every $1 \leq i, j \leq n$, and $x_{i, j}'$ such that $c_{\tA, i} = \prod_{j = 1}^n c_{\tB, j}^{x_{i, j}'}$ for every $1 \leq i, j \leq n$. 
Let $X = (x_{i, j})$ and $X' = (x_{i, j}')$. 
Both $X$ and $X'$ are invertible matrices because $X X' = I$. 

Similarly, both $\{d_{\tA, 1}, \dots, d_{\tA, m}\}$ and $\{d_{\tB, 1}, \dots, d_{\tB, m}\}$ are generating sets of $H$, there are $y_{i, j}$ such that $d_{\tB, i} = \prod_{j = 1}^m d_{\tA, j}^{y_{i, j}}$ for every $1 \leq i, j \leq m$, and $y_{i, j}'$ such that $d_{\tA, i} = \prod_{j = 1}^n d_{\tB, j}^{y_{i, j}'}$ for every $1 \leq i, j \leq m$. 
Let $Y = (y_{i, j})$ and $Y' = (y_{i, j}')$. 
Both $Y$ and $Y'$ are invertible matrices because $Y Y' = I$. 

It is easy to verify that $\tB = (X, X^t, Y') \circ \tA$. 
\end{proof}
\end{comment}


\begin{comment}
\begin{lemma}
Let $G$ be a $p$-group of nilpotent class 2 for some odd prime $p$, and $e$ be the smallest integer such that the largest order of elements in both $[G, G]$ and $G / Z(G)$ is at most $p^e$. 
Let $\tA$ and $\tB$ be two tensors derived from $G$.  $\tA$ and $\tB$ are of the same dimension and there are invertible matrices $X$ and $Y$ such that 
$\tA = (X, X^t, Y') \circ \tB$. 
\end{lemma}

\begin{proof}
By Definition~\ref{def:tensor_from_group}, the two tensors are of the same dimension. 
Let $c_{\tA, 1}, \dots, c_{\tA, n}, d_{\tA, 1}, \dots, d_{\tA, m}$ and $c_{\tB, 1}, \dots, c_{\tB, n}, d_{\tB, 1}, \dots, d_{\tB, m}$ be the elements used to obtain $\tA$ and $\tB$ following Definition~\ref{def:tensor_from_group} respectively. 

Since both $\{c_{\tA, 1}, \dots, c_{\tA, n}\}$ and $\{c_{\tB, 1}, \dots, c_{\tB, n}\}$ are generating sets of $G / Z(G)$, there are $x_{i, j}$ such that $c_{\tB, i} = \prod_{j = 1}^n c_{\tA, j}^{x_{i, j}}$ for every $1 \leq i, j \leq n$, and $x_{i, j}'$ such that $c_{\tA, i} = \prod_{j = 1}^n c_{\tB, j}^{x_{i, j}'}$ for every $1 \leq i, j \leq n$. 
Let $X = (x_{i, j})$ and $X' = (x_{i, j}')$. 
Both $X$ and $X'$ are invertible matrices because $X X' = I$. 

Similarly, both $\{d_{\tA, 1}, \dots, d_{\tA, m}\}$ and $\{d_{\tB, 1}, \dots, d_{\tB, m}\}$ are generating sets of $[G, G]$, there are $y_{i, j}$ such that $d_{\tB, i} = \prod_{j = 1}^m d_{\tA, j}^{y_{i, j}}$ for every $1 \leq i, j \leq m$, and $y_{i, j}'$ such that $d_{\tA, i} = \prod_{j = 1}^n d_{\tB, j}^{y_{i, j}'}$ for every $1 \leq i, j \leq m$. 
Let $Y = (y_{i, j})$ and $Y' = (y_{i, j}')$. 
Both $Y$ and $Y'$ are invertible matrices because $Y Y' = I$. 

One can verify that $\tB = (X, X^t, Y') \circ \tA$. 
\end{proof}
\end{comment}

%\begin{lemma}
%Let $G$ be a $p$-group of nilpotent class 2 for some odd prime $p$, $N$ be the order of $G$, and $e$ be the smallest integer such that the largest order of elements in both $[G, G]$ and $G / Z(G)$ is at most $p^e$. 
%There is an algorithm to compute a tensor derived from $G$ in polynomial time of $N$. 
%\end{lemma}

We define some notations for tensors. 
Let $\tA = (a_{i, j, k})$ be an $n \times m \times h$ tensor over $\Z / p^{e} \Z$. For $i\in[n]$, let $\mathfrak{X}_{\tA, i}$ be the $h \times m$ matrix over $\Z / p^{e} \Z$ such that $\mathfrak{X}_{\tA, i}[k, j] = a_{i, j, k}$ for any $1 \leq j \leq m, 1 \leq k \leq h$. For $j\in[m]$, let 
$\mathfrak{Y}_{\tA, j}$ be the $h \times n$ matrix over $\Z / p^{e} \Z$ such that $\mathfrak{Y}_{\tA, j}[k, i] = a_{i, j, k}$ for any $1 \leq i \leq n, 1 \leq k \leq h$. For $k\in[h]$, let $\mathfrak{Z}_{\tA, k}$ be the $n \times m$ matrix over $\Z / p^{e} \Z$ such that $\mathfrak{Z}_{\tA, k}[i, j] = a_{i, j, k}$ for any $1 \leq i \leq n, 1 \leq j \leq m$.
Let $\alpha_{\tA, 1}, \dots, \alpha_{\tA, e}$ be the integers such that the abelian group generated by $\mathfrak{Z}_{\tA, k}$ over $k\in[h]$ is isomorphic to 
\[\bigoplus_{\ell = 1}^{e} \left(\Z / p^{\ell } \Z\right)^{\alpha_{\tA, \ell}},\]
$\beta_{\tA, 1}, \dots, \beta_{\tA, e}$ be the integers such that the abelian group generated by $\mathfrak{X}_{\tA, i}$ over $i\in[n]$ is isomorphic to 
\[\bigoplus_{\ell = 1}^{e} \left(\Z / p^{\ell } \Z\right)^{\beta_{\tA, \ell}},\]
and 
$\gamma_{\tA, 1}, \dots, \gamma_{\tA, e}$ be the integers such that the abelian group generated by $\mathfrak{Y}_{\tA, j}$ over $j\in[m]$ is isomorphic to 
\[\bigoplus_{\ell = 1}^{e} \left(\Z / p^{\ell } \Z\right)^{\gamma_{\tA, \ell}}.\]
We call $\sum_{\ell = 1}^e \ell \cdot (\alpha_{\tA, \ell} + \beta_{\tA, \ell} + \gamma_{\tA, \ell})$ the \emph{complex parameter} for the tensor. 




The following theorem reduces the group isomorphism problem for class-2 nilpotent groups to the \algprobm{Alternating Matrix Space Isometry} problem. 

%Lastly, we define the following problem:
%
%\begin{quotation}
%\noindent \algprobm{Alternating Matrix Space Isometry with Stabilized Gadget}
%\noindent \textit{Input:} $\cA,\cB, \Gamma$, $\Z/p^e \Z$-submodules of $\Wedge^2(n, \Z/p^e \Z)$
%\noindent \textit{Decide:} Does there exist $P \in \GL_n(\Z/p^e Z)$ such that $P \cA P^t = \cB$ and $P \Gamma P^t = \Gamma$?
%\end{quotation}
%
%We will reduce to this problem in such a way that $\cA,\cB$ will encode our groups, and their complex parameter will be polynomial in $|G|$, while $\Gamma$ will be $O(1)$-generated, and is a ``gadget'' that enforces certain structures for the reduction.

\begin{theorem} \label{thm:GpI-to-AMSI}
The problem of testing isomorphism of odd order nilpotent groups of class 2 reduces to \algprobm{Alternating Matrix Space Isometry} over the integers modulo prime powers in polynomial time.

Moreover, if the input groups are of order $N$, then for each instance of the \algprobm{Alternating Matrix Space Isometry} problem with tensors over $p^e$ for some prime $p$ and positive integer $e$, 
then the tensors are of dimension $O(\log_p N) \times O(\log_p N) \times O(\log_p N)$ with complex parameter $O(\log_p N)$. %\jnote{I removed the base $p$ because I realized if the groups aren't $p$-groups, then there could be multiple primes.}\xnote{I change the complex parameter from linear in group order to logarithmic in group order because of the change of the complex parameter.}
\end{theorem}

Before coming to the proof, we recall a key component, which is Baer's Correspondence \cite{Bae38} (see, e.g., \cite{khukhro,naik} for expository treatments). Given a finite odd order nilpotent group $G$ of class 2, Baer's Correspondence constructs a finite Lie ring, denoted $\log(G)$, on the same underlying set as $G$, and where the Lie ring operations are defined by
\[
x + y := xy / \sqrt{[x,y]} \qquad -x := x^{-1} \qquad [x,y]_{Lie} := xyx^{-1} y^{-1},
\]
where here the expressions on the LHS are the Lie ring operations and those on the RHS are in the group. We note that $|G|$ is odd, every element has a unique square root, and since the group has class 2, $[x,y]$ is in the center, making the first expression $xy / \sqrt{[x,y]}$ well-defined. In the opposite direction, given a finite nilpotent Lie ring $L$ of odd order and class 2, Baer's Correspondence defines a group, denotes $\exp(L)$, on the same underlying set as $L$, and in which the group operations are defined by:
\[
xy := x + y + \frac{1}{2}[x,y] \qquad x^{-1} := -x,
\]
where here the LHS is in the group and the RHS is in the Lie ring $L$. Baer's Correspondence says that these two constructions are inverses of one another, and that they give an equivalence between the category of these groups and the category of these Lie rings. In particular, two groups $G,H$ are isomorphic iff $\log(G)$ and $\log(H)$ are isomorphic Lie rings. From the formulas above, it is readily seen that, in the Cayley table model, one can construct $\log(G)$ from $G$ in polynomial time, and one can construct $\exp(L)$ from $L$ in polynomial time (in fact, with the additional requirement of exponent $p$, this can also be done efficiently in the setting of matrix groups and matrix Lie rings; see, e.g., \cite[Sec.~6.2]{TI2}). Furthermore, several properties of $G$ are preserved in $\log(G)$ and vice versa:
\begin{description}
\item[Element orders.] For any integer $k$, we have $x + x + \dotsb + x$ ($k$ times) in the Lie ring is the same element as $x^k$ in the group, and thus if $x \in G$ has (multiplicative) order $p^e$, then $x \in \log(G)$ has (additive) order $p^e$. That is, in $\log(G)$, we have $p^e x = 0$ and $p^e$ is the least such integer. Thus the additive group underlying $\log(G)$ contains the information of the orders of elements of $G$.

\item[Nilpotency class.] $\log(G)$ is nilpotent of class 2.

\item[Centers.] The underlying set of $Z(G)$ is the same as the underlying set of $Z(\log(G))$. The group multiplication on $Z(G)$ is identical to the Lie ring addition operation on $Z(\log(G))$---that is, $Z(G)$ and $Z(\log(G))$ are identical as Abelian groups.

\item[Central quotients.] There is a natural isomorphism between $G / Z(G)$ and $\log(G) / Z(\log(G))$, as Abelian groups. That is, in $G/Z(G)$, the group operation is still given by multiplication in $G$, simply working modulo $Z(G)$. Similarly, since addition in $\log(G)$ is given by $x+y = xy / \sqrt{[x,y]}$ and $[x,y] \in Z(G) = Z(\log(G))$, we have that addition in $\log(G)/Z(\log(G))$ is given exactly by group multiplication, modulo $Z(G)$.
\end{description}

Given two finite Lie rings $L,L'$, an \emph{isomorphism} between them is a bijection $\varphi \colon L \to L'$ such that (1) $\varphi$ is linear, that is, $\varphi(x+y) = \varphi(x) + \varphi(y)$, (2) $\varphi([x,y]) = [\varphi(x), \varphi(y)]$ for all $x,y \in L$. As with finite groups, finite Lie rings may be given by their ``Cayley tables'', which actually consists of two tables: one for the addition operation and one for the Lie bracket.

\begin{proof}[Proof of Theorem~\ref{thm:GpI-to-AMSI}]
First, any finite nilpotent group is the direct product of its Sylow subgroups. As factoring a group into its direct products is unique (the Remak--Krull--Schmidt Theorem), and can be done in polynomial time \cite{WilsonProductsI, WilsonProductsII,KN}, for isomorphism testing we thus reduce to the case of testing isomorphism of $p$-groups of class 2, for some odd prime $p$ (which may be given as part of the input).

Next, from Baer's Correspondence, we reduce from testing isomorphism of $p$-groups of class 2 to testing isomorphism $p$-Lie rings of class 2, of the same size as the groups we started with. Let $L,L'$ denote the corresponding Lie rings that we wish to test for isomorphism.

Now, given $L$, consider the following system of linear and bilinear maps between Abelian $p$-groups. Let $U = Z(L)$, $V = L$, $W = L/Z(L)$. Let $\iota \colon U \to V$ be the natural inclusion map of $Z(L)$ as an additive subgroup of $L$, and let $\pi \colon V \to W$ be the natural quotient map $L \to L/Z(L)$ with kernel $Z(L)$. Finally, let $B \colon W \otimes W \to U$ denote the Lie bracket. Let $U',V',W',\iota', \pi', B'$ be the corresponding groups and maps correpsonding to $L'$. 

%\jnote{Wait, why are we going to all this trouble? If we are going to just say ``TI4 works over $\Z/p^e \Z$'' anyway, TI4 already lets us reduce directly from Lie algebra/ring iso to TI to AltMatSpIsom. I mean, I guess all this trouble lets us rely a \emph{little} less on just saying that TI4 works over finite local PIRs...}

\begin{lemma} \label{lem:iso-Lie-ring}
With notation as above, $L \cong L'$ if and only if there are Abelian group isomorphisms $\varphi_U \colon U \to U'$, $\varphi_V \colon V \to V'$, $\varphi_W \colon W \to W'$ such that the following conditions hold:
\begin{align} 
\varphi_V \circ \iota & = \iota' \circ \varphi_U \label{eqn:iso-condition-inj} \\
\varphi_W \circ \pi & = \pi' \circ \varphi_V \label{eqn:iso-condition-surj} \\
\varphi_U(B(x,y)) & = B'(\varphi_W(x), \varphi_W(y)) \qquad (\forall x,y \in W) \label{eqn:iso-condition-bracket}
\end{align}
\end{lemma}

This is essentially standard, but we phrase it this way as it will help us below to set up the tensors we are aiming for. Over a field, we do not need such a complicated setup, and instead can just focus on $W=L/Z(L)$ and $U=Z(L)$; by including $L$ itself in the setup, we are automatically including the Abelian group structure on $L$, which need no longer be the direct sum of $Z(L)$ and $L/Z(L)$. Conditions \eqref{eqn:iso-condition-inj} and \eqref{eqn:iso-condition-surj} then essentially say that the isomorphism on the center is just the restriction of the isomorphism on $L$, and the isomorphism on $L/Z(L)$ is the one induced from the isomorphism on $L$.

\begin{proof}[Proof of Lemma~\ref{lem:iso-Lie-ring}]
($\Rightarrow$) Suppose $\varphi \colon L \to L'$ is a Lie ring isomorphism; equivalently, $\varphi$ is an isomorphism of the additive groups such that $\varphi([x,y]) = [\varphi(x), \varphi(y)]$ for all $x,y \in L$. Since the center is a characteristic sub-Lie-ring, $\varphi(Z(L)) = Z(L')$, and therefore $\varphi$ also induces a map $L/Z(L) \to L'/Z(L')$, namely, the map which sends $x + Z(L)$ to $\varphi(x) + Z(L') = \varphi(x) + \varphi(Z(L)) = \varphi(x + Z(L))$. Let $\varphi_U = \varphi|_{Z(L)}$, $\varphi_V = \varphi$, and $\varphi_W$ be the induced map on $L/Z(L)$ just described.

Since $Z(L)$ is characteristic, $\varphi_U$ is an isomorphism of the (Abelian) Lie rings $Z(L) \cong Z(L')$, and therefore in particular an isomorphism of their underlying Abelian groups. By definition $\varphi_V = \varphi$ is an isomorphism of Lie rings, hence of Abelian groups. And then $\varphi_W$ must also be an isomorphism of Abelian groups: it is a homomorphism by construction, surjective because $\varphi$ is surjective, and injective because its kernel is $Z(L)$ by construction. 

We claim that $(\varphi_U, \varphi_V, \varphi_W)$ satisfy \eqref{eqn:iso-condition-inj}--\eqref{eqn:iso-condition-bracket}.
\begin{itemize}
\item[\ref{eqn:iso-condition-inj}: ] Suppose $u \in U = Z(L)$. Then we have $\varphi_V(\iota(u)) = \varphi(u)$ and $\iota'(\varphi_U(u)) = \iota'(\varphi(u)) = \varphi(u)$, where here we have identified elements of $U = Z(L)$ with the corresponding elements of $L$, so that e.g. $\iota(u) = u$ and $\iota'(u') = u'$ for $u \in Z(L), u' \in Z(L')$.

\item[\ref{eqn:iso-condition-surj}: ] Let $x \in L$. Then we have $\pi(x) = x + Z(L)$ (considered as an element of $L/Z(L)$). So we have 
\begin{align*}
\varphi_W(\pi(x)) & = \varphi_W(x + Z(L)) = \varphi(x + Z(L)) = \varphi(x) + \varphi(Z(L)) \\ & = \varphi(x) + Z(L') = \pi'(\varphi(x)) = \pi'(\varphi_V(x)).
\end{align*}

\item[\ref{eqn:iso-condition-bracket}: ] Let $x,y \in L$.  Then we have
\begin{align*}
\varphi_U(B(x + Z(L),y + Z(L))) & = \varphi_U([x,y]) & \text{(since $[Z(L),*]=[*,Z(L)]=0$)} \\
 & = \varphi([x,y]) & \text{(since $[x,y] \in Z(L)$)} \\
 & = [\varphi(x), \varphi(y)] & \text{(since $\varphi$ is a homomorphism)}\\
 & = [\varphi(x) + Z(L'), \varphi(y) + Z(L')] \\
 &  = B'(\varphi_W(x), \varphi_W(y)).
\end{align*}

\end{itemize}

($\Leftarrow$) Conversely, suppose $(\varphi_U, \varphi_V, \varphi_W)$ satisfy \eqref{eqn:iso-condition-inj}--\eqref{eqn:iso-condition-bracket}. Let $\varphi = \varphi_V$; we claim that $\varphi$ is in fact an isomorphism of Lie rings. By assumption, $\varphi_V = \varphi$ is an isomorphism of Abelian groups, so all that remains to show is that $\varphi$ preserves the Lie bracket, i.e., that $\varphi([x,y]) = [\varphi(x), \varphi(y)]$. For all $x,y \in L$, we have:
\begin{align*}
\varphi([x,y]) & = \varphi_V(\iota(B(x + Z(L),y + Z(L)))) & \text{(definition of $B, \varphi, \iota$)}\\ 
 & = \iota'(\varphi_U(B(x + Z(L),y + Z(L)))) & \text{(by \eqref{eqn:iso-condition-inj})} \\
 & = \iota'(B'(\varphi_W(\pi(x)), \varphi_W(\pi(y)))) & \text{(by \eqref{eqn:iso-condition-bracket})} \\
 & = [\varphi_W(\pi(x)), \varphi_W(\pi(y))] & \text{(by definition of $B',\iota'$)} \\
 & = [\pi'(\varphi(x)), \pi'(\varphi(y))] & \text{(by \eqref{eqn:iso-condition-surj} and def'n of $\varphi$)} \\
 & = [\varphi(x) + Z(L'), \varphi(y) + Z(L')] & \text{(by definition of $\pi'$)} \\
 & = [\varphi(x), \varphi(y)] & \text{(since $[Z(L'),*] = [*, Z(L')] = 0$.)}.
\end{align*}

This completes the proof of Lemma~\ref{lem:iso-Lie-ring}.
\end{proof}


\begin{lemma} \label{lem:Lie-ring-to-TI}
There is a polynomial-time algorithm, which, given a nilpotent $p$-Lie ring $L$ of class 2 and exponent $p^e$ (by Cayley tables or by basis and structure coefficients), produces a tensor $T_L$ in $(L \oplus L \oplus L/Z(L) \oplus L/Z(L) \oplus \Z/p^e) \otimes (L \oplus L \oplus L/Z(L) \oplus L/Z(L) \oplus \Z/p^e)  \otimes (\Z/p^e \oplus Z(L) \oplus Z(L))$ such that $L \cong L'$ if and only if there are $X_1 \in \aut(Z(L))$, $X_2 \in \aut(L_+)$ (automorphism of the underlying Abelian group of $L$), and $X_3 \in \aut(L/Z(L))$ such that
\begin{equation} \label{eqn:Lie-ring-to-TI}
(X_2 \oplus X_2^{-1} \oplus X_3^{-1} \oplus X_3 \oplus 1, (X_2 \oplus X_2^{-1} \oplus X_3^{-1} \oplus X_3 \oplus 1)^t, 1 \oplus X_1 \oplus X_1^{-1}) \circ T_L = T_{L'}.
\end{equation}
Furthermore, each frontal slice of the tensors $T_L, T_{L'}$ is an alternating matrix.
\end{lemma}

\begin{proof}
View $T_L$ as a block tensor, as indicated by the shape $(L \oplus L \oplus L/Z(L) \oplus L/Z(L) \oplus \Z/p^e) \otimes (L \oplus L \oplus L/Z(L) \oplus L/Z(L) \oplus \Z/p^e)  \otimes (\Z/p^e \oplus Z(L) \oplus Z(L))$; that is, there are $75$ blocks arranged in a $5 \times 5 \times 3$ fashion, and e.g. the (1,1,1) block is $L \otimes L \otimes \Z/p^e$, and the (3,3,2) block is $L/Z(L) \otimes L/Z(L) \otimes Z(L)$, and the (5,2,3) block is $\Z/p^e \otimes L \otimes Z(L)$. Throughout, choose a basis for $L$, a basis of $Z(L)$, and a basis of $L/Z(L)$. 
\begin{itemize}
\item Let $\tB \colon L / Z(L) \otimes L / Z(L) \to Z(L)$ be the commutator map of $L$; we place $\tB$, written in the chosen bases, in the (3,3,2) block. 

\item Let $M_\pi$ be the matrix corresponding to the projection map $\pi \colon L \to L/Z(L)$ (again, in the chosen bases); we place $M_{\pi}$ in the (4,2,1) block, and $-M_{\pi}^t$ in the (2,4,1) block. 

\item Let $M_{\iota}$ be the matrix corresponding to the inclusion map $\iota \colon Z(L) \to L$; we place $M_{\iota}$ in the (5,1,3) block and $-M_{\iota}^t$ in the (1,5,3) block.
\end{itemize}

Now we show that this tensor $T_L$ has the desired property, namely that $L \cong L'$ if and only if $T_L$ and $T_{L'}$ are isomorphic by a transformation as in the statement of the lemma.

($\Rightarrow$) Suppose $L \cong L'$. By Lemma~\ref{lem:iso-Lie-ring}, there are isomorphisms $\varphi_U \colon Z(L) \to Z(L')$, $\varphi_V \colon L \to L'$, and $\varphi_W \colon L/Z(L) \to L' / Z(L')$ satisfying equations \eqref{eqn:iso-condition-inj}--\eqref{eqn:iso-condition-bracket}. Let $X_1 = \varphi_U$, $X_2 = \varphi_V$, and $X_3 = \varphi_W$. We claim that $X_1, X_2, X_3$ satisfy \eqref{eqn:Lie-ring-to-TI}.
\begin{itemize}
\item Consider the action on $\tB$ in the (3,3,2) block of $T_L$. This is acted on by $\tB \mapsto X_3^{-1} \tB^{X_1} X_3^{-T}$, where here the exponent $\tB^{X_1}$ means that $X_1$ is used to take linear combinations of frontal slices of $\tB$. In terms of the bilinear form, this has the effect of replacing $B(x,y)$ by $X_1 B(X_3^{-1} x, X_3^{-1} y) = \varphi_U(B(\varphi_W^{-1}(x), \varphi_W^{-1}(y))$. By \eqref{eqn:iso-condition-bracket}, this is equal to $B'(x,y)$, which agrees with the (3,3,2) block of $T_{L'}$.

\item Consider the action on $M_{\pi}$ in the (1,4,1) block.  This transforms into $M_\pi \mapsto X_3 M_\pi X_2^{-T}$. But this is the same as $\varphi_W \circ \pi \circ \varphi_V^{-1}$, which by \eqref{eqn:iso-condition-surj} is equal to $\pi'$, hence making the (1,4,1) block into $M_{\pi'}$, the (1,4,1) block of $T_{L'}$. In the (4,1,1) block, the same equation occurs just with everything transposed and negated, resulting in $-M_{\pi'}$, again agreeing with the (4,1,1) block of $T_{L'}$.

\item Similarly, in the (5,1,3) block, we have $M_{\iota} \mapsto X_2 M_{\iota} X_1^{-1}$, corresponding to $\varphi_V \circ \iota \circ \varphi_U^{-1}$, and by \eqref{eqn:iso-condition-inj}, this is equal to $\iota'$, hence the resulting matrix equals $M_{\iota'}$, which is precisely the (5,1,3) block of $T_{L'}$. As in the previous case, the (1,5,3) block is simply the negative transpose in both $T_L$ and $T_{L'}$.
\end{itemize}

($\Leftarrow$) Given invertible $X_1, X_2, X_3$ satisfying \eqref{eqn:Lie-ring-to-TI}, define $\varphi_U = X_1, \varphi_V = X_2, \varphi_W = X_3$. Reading the preceding argument in reverse shows that $\varphi_U, \varphi_V, \varphi_W$ satisfy \eqref{eqn:iso-condition-inj}--\eqref{eqn:iso-condition-bracket}, and therefore by Lemma~\ref{lem:iso-Lie-ring}, we have that $L \cong L'$. This completes the proof of Lemma~\ref{lem:Lie-ring-to-TI}.
\end{proof}

Deciding if two tensors as in Lemma~\ref{lem:Lie-ring-to-TI} are equivalent under the action specified there is an instance of \algprobm{Partitioned \& Linked Alternating Matrix Space Isometry} over $\Z/p^e \Z$, but where the tensors need not be homocyclic (the modules need not be free). The latter problem is the same as \algprobm{Alternating Matrix Space Isometry} except that the matrices doing the isometry must have a block form (specified as part of the input; this is why it is called ``partitioned''), and certain blocks of the matrices doing the isometry are required to be equal, inverses, or transposes of one another (why it is called ``linked''). In \cite{TI4}, over fields gadgets were presented that reduces the latter problem to \algprobm{Alternating Matrix Space Isometry}. We claim that those same gadgets work in our setting, and the only key adaptation we need is Lemma~\ref{lem:invertible} and Remark~\ref{rem:TI4}, which allows us to reduce to (ordinary, non-partitioned, non-linked, but still not necessarily homocyclic)  \algprobm{Alternating Matrix Space Isometry} over $\Z/p^e \Z$. 

Finally, we reduce to \algprobm{Alternating Matrix Space Isometry} over free (equivalently, homocyclic) $\Z/p^e \Z$-modules. 

\begin{definition}\label{def:tensor_from_group}
Let $G$ and $H$ be two Abelian $p$-groups for some odd prime $p$, 
$f: G \times G \rightarrow H$ be a bilinear map, and $\tA$ be the tensor in $G \otimes G\otimes H$ with respect to $f_\tA$. 
Let $e$ be an integer such that the largest order of elements in both $G$ and $H$ is at most $p^e$. 
Suppose $G$ is isomorphic to 
\[\bigoplus_{\ell = 1}^{e} \left(\Z / p^{\ell } \Z\right)^{\alpha_\ell},\]
and $H$ is isomorphic to 
\[\bigoplus_{\ell = 1}^{e} \left(\Z / p^{\ell } \Z\right)^{\beta_\ell}.\]
Let $n = \sum_{\ell = 1}^{e} \alpha_\ell$ and $m = \sum_{\ell = 1}^{e} \beta_\ell $. 
An $n \times n \times m$ tensor $\tB = (b_{i, j, k})$ over $\Z / p^e \Z$ is a tensor derived from $\tA$ if there exist
$c_1, \dots, c_n \in G$ and $d_1, \dots, d_m \in H$ satisfying the following conditions:
\begin{enumerate}
\item $c_1, \dots, c_n$ generate $G$ and for each $1 \leq \ell \leq e$, $c_{j}$ is of order $p^\ell$ for any $(\sum_{i = 1}^{\ell - 1} \alpha_i) + 1 \leq j \leq \sum_{i = 1}^{\ell} \alpha_i$. (Equivalently, the $c_i$ are a basis for $G$ viewed as a $\Z/p^e \Z$-module.)
\item $d_1, \dots, d_m$ generate $H$ and for each $1 \leq \ell \leq e$, $d_{j}$ is of order $p^\ell$ for any $(\sum_{i = 1}^{\ell - 1} \beta_i) + 1 \leq j \leq \sum_{i = 1}^{\ell} \beta_i$. (Equivalently, the $d_i$ are a basis for $H$ viewed as a $\Z/p^e \Z$-module.)
\item For any $1 \leq i, j \leq n, 1 \leq k \leq m$, $b_{i, j, k} = f_\tA(c_i, c_j, d_k)\cdot p^{e - t_k}$, where $t_k$ is the integer such that the order of $d_k$ is $p^{t_k}$.%\ynote{I am a bit confused here: why the padded $p$ power is in terms of $d_k$, but not relying on $c_j$?}
\end{enumerate}



%Suppose $[G, G]$ is isomorphic to 
%\[\bigoplus_{\ell = 1}^{e} \left(\Z / p^{\ell } \Z\right)^{\alpha_\ell},\]
%and $G / Z(G)$ is isomorphic to 
%\[\bigoplus_{\ell = 1}^{e} \left(\Z / p^{\ell } \Z\right)^{\beta_\ell}.\]
%Let $m = \sum_{\ell = 1}^{e} \alpha_\ell$ and $n = \sum_{\ell = 1}^{e} \beta_\ell $. 



%and $e$ be the smallest integer such that the largest order of elements in both $[G, G]$ and $G / Z(G)$ is at most $p^e$. 
%Suppose $[G, G]$ is isomorphic to 
%\[\bigoplus_{\ell = 1}^{e} \left(\Z / p^{\ell } \Z\right)^{\alpha_\ell},\]
%and $G / Z(G)$ is isomorphic to 
%\[\bigoplus_{\ell = 1}^{e} \left(\Z / p^{\ell } \Z\right)^{\beta_\ell}.\]
%Let $m = \sum_{\ell = 1}^{e} \alpha_\ell$ and $n = \sum_{\ell = 1}^{e} \beta_\ell $. 

%An $n \times n \times m$ tensor $\tA = (a_{i, j, k})$ over $\Z / p^e \Z$ is a tensor derived from $G$ if there exist
%$c_1, \dots, c_n \in G / Z(G)$ and $d_1, \dots, d_m$ satisfying the following conditions:
%\begin{enumerate}
%\item $c_1, \dots, c_n$ generate $G / Z(G)$ and for each $1 \leq \ell \leq e$, $c_{j}$ is of order $p^\ell$ for any $(\sum_{i = 1}^{\ell - 1} \beta_i) + 1 \leq j \leq \sum_{i = 1}^{\ell} \beta_i$. 
%\item $d_1, \dots, d_m$ generate $[G, G]$ and for each $1 \leq \ell \leq e$, $d_{j}$ is of order $p^\ell$ for any $(\sum_{i = 1}^{\ell - 1} \alpha_i) + 1 \leq j \leq \sum_{i = 1}^{\ell} \alpha_i$. 
%\item For any $1 \leq i, j \leq n$, $[c_i, c_j] = \prod_{k = 1}^md_k^{a_{i, j, k}} \cdot p^{e - t}$, where $t$ is the integer such that the order of $d_k$ is $p^t$. 
%\end{enumerate}
\end{definition}



\begin{definition}\label{def:matrix_from_group}
Let $G$ and $H$ be two Abelian $p$-groups for some odd prime $p$, 
$f: G \rightarrow H$ be a linear map, and  
$A$ be a matrix in $G \otimes H$ with respect to $f$. 
Let $e$ be an integer such that the largest order of elements in both $G$ and $H$ is at most $p^e$. 
Suppose $G$ is isomorphic to 
\[\bigoplus_{\ell = 1}^{e} \left(\Z / p^{\ell } \Z\right)^{\alpha_\ell},\]
and $H$ is isomorphic to 
\[\bigoplus_{\ell = 1}^{e} \left(\Z / p^{\ell } \Z\right)^{\beta_\ell}.\]
Let $n = \sum_{\ell = 1}^{e} \alpha_\ell$ and $m = \sum_{\ell = 1}^{e} \beta_\ell $. 
An $n \times m$ matrix  $B = (b_{i, j})$ over $\Z / p^e \Z$ is a matrix derived from $A$ if there exist
$c_1, \dots, c_n \in G$ and $d_1, \dots, d_m \in H$ satisfying the following conditions:
\begin{enumerate}
\item $c_1, \dots, c_n$ generate $G$ and for each $1 \leq \ell \leq e$, $c_{j}$ is of order $p^\ell$ for any $(\sum_{i = 1}^{\ell - 1} \alpha_i) + 1 \leq j \leq \sum_{i = 1}^{\ell} \alpha_i$. 
\item $d_1, \dots, d_m$ generate $H$ and for each $1 \leq \ell \leq e$, $d_{j}$ is of order $p^\ell$ for any $(\sum_{i = 1}^{\ell - 1} \beta_i) + 1 \leq j \leq \sum_{i = 1}^{\ell} \beta_i$. 
\item For any $1 \leq i, j \leq n, 1 \leq k \leq m$, $b_{i, j} = f_A(c_i, d_j)\cdot p^{e - t}$, where $t$ is the integer such that the order of $d_j$ is $p^t$. 
\end{enumerate}



%Suppose $[G, G]$ is isomorphic to 
%\[\bigoplus_{\ell = 1}^{e} \left(\Z / p^{\ell } \Z\right)^{\alpha_\ell},\]
%and $G / Z(G)$ is isomorphic to 
%\[\bigoplus_{\ell = 1}^{e} \left(\Z / p^{\ell } \Z\right)^{\beta_\ell}.\]
%Let $m = \sum_{\ell = 1}^{e} \alpha_\ell$ and $n = \sum_{\ell = 1}^{e} \beta_\ell $. 



%and $e$ be the smallest integer such that the largest order of elements in both $[G, G]$ and $G / Z(G)$ is at most $p^e$. 
%Suppose $[G, G]$ is isomorphic to 
%\[\bigoplus_{\ell = 1}^{e} \left(\Z / p^{\ell } \Z\right)^{\alpha_\ell},\]
%and $G / Z(G)$ is isomorphic to 
%\[\bigoplus_{\ell = 1}^{e} \left(\Z / p^{\ell } \Z\right)^{\beta_\ell}.\]
%Let $m = \sum_{\ell = 1}^{e} \alpha_\ell$ and $n = \sum_{\ell = 1}^{e} \beta_\ell $. 

%An $n \times n \times m$ tensor $\tA = (a_{i, j, k})$ over $\Z / p^e \Z$ is a tensor derived from $G$ if there exist
%$c_1, \dots, c_n \in G / Z(G)$ and $d_1, \dots, d_m$ satisfying the following conditions:
%\begin{enumerate}
%\item $c_1, \dots, c_n$ generate $G / Z(G)$ and for each $1 \leq \ell \leq e$, $c_{j}$ is of order $p^\ell$ for any $(\sum_{i = 1}^{\ell - 1} \beta_i) + 1 \leq j \leq \sum_{i = 1}^{\ell} \beta_i$. 
%\item $d_1, \dots, d_m$ generate $[G, G]$ and for each $1 \leq \ell \leq e$, $d_{j}$ is of order $p^\ell$ for any $(\sum_{i = 1}^{\ell - 1} \alpha_i) + 1 \leq j \leq \sum_{i = 1}^{\ell} \alpha_i$. 
%\item For any $1 \leq i, j \leq n$, $[c_i, c_j] = \prod_{k = 1}^md_k^{a_{i, j, k}} \cdot p^{e - t}$, where $t$ is the integer such that the order of $d_k$ is $p^t$. 
%\end{enumerate}
\end{definition}


\begin{lemma}\label{lem:reduction_to_homocyclic}
Let $G, H, G', H'$ be four Abelian $p$-groups such that $G \cong G'$ and $H \cong H'$. Then we have 
\begin{enumerate}
\item Let $\tA, \tA'$ be two tensors in $G \otimes G  \otimes H$ and $G' \otimes G' \otimes H'$, and 
$\tB$ and $\tB'$ be two tensors over $\Z / p^e \Z$ for some integer $e$ derived from $\tA$ and $\tA'$ respectively. 
Then there are $X_G \in \aut(G)$ and $Y_H \in \aut(H)$ such that 
$\tA' = (X_G, X_G^t, Y_H) \circ \tA$ 
%$\tA$ and $\tA'$ are isomorphic under actions in the automorphisms of $G$ and $H$ 
if and only if there are invertible matrices $X \in \GL(n, \Z / p^e \Z)$ and $Y \in \GL(m, \Z / p^e \Z)$ such that $\tB' = (X, X^t, Y) \circ \tB$, where $n$ is the size of a basis for $G$ and $m$ is the size of a basis for $H$. 
%their derived tensors are isomorphic under the actions in $\GL(n, \Z / p^3 \Z) \times \GL(n, \Z / p^3 \Z)$.
\item Let $A, A'$ be two matrices in $G  \otimes H$ and $G' \otimes H'$, and 
$B$ and $B'$ be two matrices over $\Z / p^e \Z$ for some integer $e$ derived from $B$ and $B'$ respectively. 
Then there are $X_G \in \aut(G)$ and $Y_H \in \aut(H)$ such that 
$A' = X_G A Y_H$ %~\xnote{I do not know if the notations used here is fine.}
%$\tA$ and $\tA'$ are isomorphic under actions in the automorphisms of $G$ and $H$ 
if and only if there are invertible matrices $X \in \GL(n, \Z / p^e \Z)$ and $Y \in \GL(m, \Z / p^e \Z)$ such that $B'= X  B Y$, where $n$ is the size of a basis for $G$ and $m$ is the size of a basis for $H$. 
%their derived tensors are isomorphic under the actions in $\GL(n, \Z / p^3 \Z) \times \GL(n, \Z / p^3 \Z)$.

\end{enumerate}
\end{lemma}
\begin{proof}
%Let $\tB$ and $\tB'$ are tensors derived 
By Definition~\ref{def:tensor_from_group}, the $\tB$ and $\tB'$ are of dimension $n \times n \times m$. 
Let $c_{\tA, 1}, \dots, c_{\tA, n}, d_{\tA, 1}, \dots, d_{\tA, m}$ and $c_{\tA', 1}, \dots, c_{\tA', n}, d_{\tA', 1}, \dots, d_{\tA', m}$ be the elements used to obtain $\tB$ and $\tB'$  following Definition~\ref{def:tensor_from_group} respectively. 

%\ynote{The proof below seems to miss some argument for why the padded $p$ powers as in Definition~\ref{def:tensor_from_group} do not ``matter.''}
We first show that if $\tA$ and $\tA'$ are isomorphic, then there are $X$ and $Y$ such that $\tB' = (X, X^t, Y)\circ \tB$. 
Since $\tA$ and $\tA'$ are isomorphic, %there are automorphism $A_G \in \aut(G)$ and $A_H \in \aut(H)$ such that $(A_G, A_G^t, A_H) \circ \tA = \tA'$. 
%Hence, 
 there are $x_{i, j} \in \Z / p^e \Z$ such that $X_G(c_{\tA, i}) = \sum_{j = 1}^n {x_{i, j}} c_{\tA, j}$ %\ynote{One of $c_{\tA, i}$ and $c_{\tA, j}$ should have $\tA'$ in the subscript?} 
 for every $1 \leq i \leq n$. %, and $x_{i, j}'$ such that $c_{\tA, i} = \prod_{j = 1}^n c_{\tB, j}^{x_{i, j}'}$ for every $1 \leq i, j \leq n$. 
 Let $X = (x_{i, j})$. % and $X' = (x_{i, j}')$. 
 $X$ is in $\GL(n, \Z / p^e \Z)$. % because 
%Both $X$ and $X'$ are invertible matrices because $X X' = I$. 
Similarly, there are $y_{i, j} \in \Z / p^e \Z$ such that $Y_H(d_{\tA, i}) = \sum_{j = 1}^m y_{i, j} d_{\tA, j}$ for every $1 \leq i \leq m$. %, and $x_{i, j}'$ such that $c_{\tA, i} = \prod_{j = 1}^n c_{\tB, j}^{x_{i, j}'}$ for every $1 \leq i, j \leq n$. 
 Let $Y_0 = (y_{i, j})$. % and $X' = (x_{i, j}')$. 
 $Y_0$ is in $\GL(n, \Z / p^e \Z)$. % because 
%Similarly, both $\{d_{\tA, 1}, \dots, d_{\tA, m}\}$ and $\{d_{\tB, 1}, \dots, d_{\tB, m}\}$ are generating sets of $H$, there are $y_{i, j}$ such that $d_{\tB, i} = \prod_{j = 1}^m d_{\tA, j}^{y_{i, j}}$ for every $1 \leq i, j \leq m$, and $y_{i, j}'$ such that $d_{\tA, i} = \prod_{j = 1}^n d_{\tB, j}^{y_{i, j}'}$ for every $1 \leq i, j \leq m$. 
%Let $Y = (y_{i, j})$ and $Y' = (y_{i, j}')$. 
%Both $Y$ and $Y'$ are invertible matrices because $Y Y' = I$. 
It is easy to verify that $\tB' = (X, X^t, Y_0^{-1}) \circ \tB$. 

Now we show that if there are invertible $X$ and $Y$ such that $\tB' = (X, X^t, Y)\circ \tB$, then $\tA$ and $\tA'$ are isomorphic. 
Denote $X = (x_{i, j})$ and $Y' = Y^{-1} = (y_{i, j}')$. 
Let $X_G$ be the linear function for $G$ such that $X_G(c_{\tA, i}) = \sum_{j = 1}^{n} x_{i, j} c_{\tA, j}$, and 
$Y_H$ be the linear function for $H$ such that $Y_H(d_{\tA, i}) = \sum_{j = 1}^{m} y'_{i, j} d_{\tA, j}$. 
One can verify that both $X_G$ and $Y_H$ are automorphisms for $G$ and $H$ respectively, and $\tA' = (X_G, X_G^t, Y_H)$. 

The proof for the second case of the lemma is similar to the first case. 
This completes the proof of Lemma~\ref{lem:reduction_to_homocyclic}.
\end{proof}

\begin{definition}
Let $G_1, \dots, G_k, H_1, \dots, H_\ell$ be a sequence of Abelian $p$-groups. 
Let $\tA$ be a tensor in %\ynote{Was: over} 
$(G_1 \oplus \dots \oplus G_k) \otimes (G_1 \oplus \dots \oplus G_k)  \otimes (H_1 \oplus \dots \oplus H_\ell)$ such that each non-zero subtensor on $G_i \otimes G_{j} \otimes H_u$ satisfies one of the following conditions:
\begin{enumerate}
\item $i = i'$ and the subtensor encodes a bilinear map. 
\item One of $G_i, G_j, H_u$ is isomorphic to 1, and the subtensor encodes a linear map. 
\end{enumerate}
The derived tensor over $\Z / p^e \Z$ for $\tA$ is obtained by replacing each non-zero sub-tensor by its derived tensor. 
\end{definition}

\begin{corollary}
Let $G_1, \dots, G_k, H_1, \dots, H_\ell$ and $G_1', \dots, G_k', H_1', \dots, H_\ell'$ be two sequences of Abelian $p$-groups such that $G_i \cong G_i'$ for every $1\leq i\leq k$ and $H_j \cong H_j'$ for every $1 \leq j \leq \ell$. 
Let $\tA$ and $\tA'$ be tensors in $(G_1 \oplus \dots \oplus G_k) \otimes (G_1 \oplus \dots \oplus G_k)  \otimes (H_1 \oplus \dots \oplus H_\ell)$ and  $(G_1' \oplus \dots \oplus G_k') \otimes (G_1' \oplus \dots \oplus G_k')  \otimes (H_1' \oplus \dots \oplus H_\ell')$ such that for each non-zero subtensor on $G_i \otimes G_{j} \otimes H_u$ (or $G_i' \otimes G_{j}' \otimes H_u'$) satisfying one of the following conditions:
\begin{enumerate}
\item $i = i'$ and the subtensor encodes a bilinear map. 
\item One of $G_i, G_j, H_u$ (or $G_i', G_j', H_u'$) is isomorphic to 1, and the subtensor encodes a linear map. 
\end{enumerate}
Then the two tensors are isomorphic if and only if their derived tensors over $\Z / p^e \Z$ are isomorphic.  
\end{corollary}

The dimension of the tensor obtained and the complex parameter is obtained by the construction of the tensor.%\ynote{Some estimation on the parameter bounds would be needed here?}

This completes the proof of Theorem~\ref{thm:GpI-to-AMSI}.
\end{proof}



% !TEX root = TI5.tex
\subsection{Reduction to matrix tuple isometry}\label{sec:tensor_iso_to_matrix_tuple}
\begin{comment}
We define some notations for tensors. 
Let $\tA = (a_{i, j, k})$ be an $n \times m \times h$ tensor over $\Z / p^{e} Z$. Let $\mathfrak{X}_{\tA, i}$ be the $h \times m$ matrix over $\Z / p^{e} \Z$ such that $\mathfrak{X}_{\tA, i}[k, j] = a_{i, j, k}$ for any $1 \leq j \leq m, 1 \leq k \leq h$, 
$\mathfrak{Y}_{\tA, j}$ be the $h \times n$ matrix over $\Z / p^{e} \Z$ such that $\mathfrak{Y}_{\tA, j}[k, i] = a_{i, j, k}$ for any $1 \leq i \leq n, 1 \leq k \leq h$, 
and $\mathfrak{Z}_{\tA, k}$ be the $n \times m$ matrix over $\Z / p^{e} \Z$ such that $\mathfrak{Z}_{\tA, k}[i, j] = a_{i, j, k}$ for any $1 \leq i \leq n, 1 \leq j \leq m$.
Let $\alpha_{\tA, 1}, \dots, \alpha_{\tA, e}$ be the integers such that the abelian group generated by $\mathfrak{Z}_{\tA, k}$ for all the $k$ is isomorphic to 
\[\bigoplus_{\ell = 1}^{e} \left(\Z / p^{\ell } \Z\right)^{\alpha_{\tA, \ell}},\]
$\beta_{\tA, 1}, \dots, \beta_{\tA, e}$ be the integers such that the abelian group generated by $\mathfrak{X}_{\tA, i}$ for all the $i$ is isomorphic to 
\[\bigoplus_{\ell = 1}^{e} \left(\Z / p^{\ell } \Z\right)^{\beta_{\tA, \ell}},\]
and 
$\gamma_{\tA, 1}, \dots, \gamma_{\tA, e}$ be the integers such that the abelian group generated by $\mathfrak{Y}_{\tA, j}$ for all the $j$ is isomorphic to 
\[\bigoplus_{\ell = 1}^{e} \left(\Z / p^{\ell } \Z\right)^{\gamma_{\tA, \ell}},\]

We call $p^{\sum_{\ell = 1}^e \ell \cdot (\alpha_{\tA, \ell} + \beta_{\tA, \ell} + \gamma_{\tA, \ell})}$ the complex parameter for the tensor. 
\end{comment}
In this section, we prove the following theorem. 



\begin{theorem}\label{thm:tensor_iso_algo}
Let $\tA, \tB$ be two $n \times n \times m$ alternating tensors over $\Z / p^e \Z$ with the same complex parameter $\mu$. 
There is an algorithm to determine if 
there are invertible matrices $P \in \GL(n, \Z / p^e \Z)$ and $Q \in \GL(m, \Z / p^e \Z)$ such that 
$\tB = (P, P, Q) \circ \tA$ in time 
\[p^{O(\mu^{17 / 9} \cdot (\log p)^{8/9})} \cdot t,\]
where $t$ is the running time to determine whether two alternating matrix tuples $(A_1, \dots, A_k)$ and $(B_1, \dots, B_k)$ in $\M(n', \Z / p^e \Z)^{k}$ (for some $n'$ and $k$ that are polynomial in $n$ and $m$) has an invertible matrix $P \in \GL(n', \Z / p^e \Z)$ such that $A_i = P B_i P^t$ for any $1 \leq i \leq k$.
\end{theorem}

\subsubsection{Individualizing tensors over $\Z/p^e \Z$}

In this section, we consider the following problem: Let $\mathfrak{A}$ be a $\Z / p^e \Z$-submodule of $\M(n\times m, \Z / p^e \Z)$, and $r$ be an integer at most $\min\{n, m\}$. We want to choose two matrices $L$ and $R$ such that for any matrix $A \in \mathfrak{A}$ with inner rank at least $r$, $LAR$ is a non-zero matrix. 
Our goal is to choose $L$ and $R$ such that the number of rows of $L$ and the number of columns of $R$ is small. 

\begin{lemma}\label{lem:sampling_p_to_e}
Let $A$ be a matrix in $\M(m\times n,  \Z / p^e \Z)$ with inner rank $r$ and $L$ be a matrix in $M(t\times m,  \Z / p^e \Z)$ such that $LA$ is of inner rank $r'$ for some integer $r' < r$. 
For a randomly sampled row vector $q \in (\Z / p^e \Z)^m$, 
\[\begin{bmatrix}L \\ q\end{bmatrix}A\] is of inner rank $r' + 1$ with probability at least $1 - (1/p)^{r - r'}$. 
\end{lemma}
\begin{proof}
\begin{comment}
We first consider the case of $r' = r - 1$. 
By the definition of inner rank, there is a  submodule $D$ of $(\Z / p^e \Z)^n$ generated by $r'$ elements containing the row-span of $LA$. 
Let $S_{D} = \{q \in (\Z / p^e \Z)^m : q A \in D\}$. 
By the linearity, $S_D$ corresponds to a submodule of $(\Z / p^e \Z)^m$. 

$S_D$ is not equal to $ (\Z / p^e \Z)^m$ because otherwise $P$ is of inner rank $r$. 
Let $q$ be an arbitrary element in $ (\Z / p^e \Z)^m \setminus S_D$. 
Clearly, \[\begin{bmatrix}L \\ q\end{bmatrix}A\] is a matrix of inner rank $r$. 
In addition, for any $i \in \{1, \dots, p - 1\}$, \[\begin{bmatrix}L \\ i\cdot q\end{bmatrix}A\] is also of inner rank $r$. 
Hence, we have
\[\frac{|(\Z / p^e \Z)^m|}{|S_D|} \geq p.\]
So, for a randomly sampled vector $q \in (\Z / p^e \Z)^m$, 
\[\begin{bmatrix}L \\ q\end{bmatrix}A\] is a matrix of inner rank $r' + 1$ with probability at least $1 - p$. 
\end{comment}

%Now consider the case of $r' < r - 1$. 
Let $L_{r'}$ be a matrix satisfying the following conditions:
\begin{enumerate}
\item $\rank(L_{r'} A) =\rank(L A) = r'$ and   $\rowspan(L A) \subset \rowspan(L_{r'} A)$, where $\rowspan(A)$ denotes the $\Z / p^e \Z$-span for the rows of matrix $A$.
\item There is no row vector $q \in (\Z / p^e \Z)^m$  such that 
\[\rank\left(\begin{bmatrix} L_{r'} \\ q \end{bmatrix} A\right) = r' \text{ and } 
\rowspan(L_{r'} A) \subsetneq \rowspan\left(\begin{bmatrix} L_{r'} \\ q \end{bmatrix} A\right).\]
\end{enumerate}
For all $r' + 1 \leq r'' \leq r$, let
\[L_{r''}  = \begin{bmatrix}L_{r'' - 1} \\ L_{r''}' \end{bmatrix}\]
where $L_{r''}'$ is a matrix with row vectors in $(\Z / p^e \Z)^m$ 
such that the following conditions hold: 
\begin{enumerate}
\item $\rank(L_{r''} A) = r''$. 
\item There is no $q_{r''}'$ as a row vector in $(\Z / p^e \Z)^m$  such that 
$\rowspan(L_{r''} A) \subsetneq \langle \rowspan(L_{r''}) \cup \{q_{r''}'\} \rangle$.
\item There is no row vector $q \in (\Z / p^e \Z)^m$  such that 
\[\rank\left(\begin{bmatrix} L_{r''} \\ q \end{bmatrix} A\right) = r'' \text{ and } 
\rowspan(L_{r''} A) \subsetneq \rowspan\left(\begin{bmatrix} L_{r''} \\ q \end{bmatrix} A\right).\]

\end{enumerate}
%By induction that $L_{r'' - 1} A$ is of inner rank $r'' - 1$ and by the fact that $A$ is of rank $r$, $q_{r''}$ exists for every $r' \leq r'' \leq r$. 
By the fact that  $A$ is of rank $r$, such $L_{r'}, L_{r' + 1}, \dots, L_{r}$ exist. 

%a matrix containing all the rows of $L_{r'' - 1}$ satisfying that $L_{r''} A$ is of inner rank $r''$.
%Let $D_{r'}$ be an arbitrary submodule of $(\Z / p^e \Z)^n$ generated by $r'$ elements containing row-span of $L_{r'}P$,
%and $D_{r''}$ all $r' + 1 \leq r'' \leq r$ be an arbitrary submodule of $(\Z / p^e \Z)^n$ generated by $r''$ elements containing both row-span of $L_{r''}P$ and $D_{r'' - 1}$. 
Let $S_{r''} = \{q \in (\Z / p^e \Z)^m : q\cdot A \in \rowspan(L_{r''} A) \}$ for all $r' \leq r'' \leq r$.
For any $r' \leq r'' \leq r - 1$, any vector $q \in (\Z / p^e \Z)^m \setminus S_{r''}$ satisfies that the inner rank of 
\[\begin{bmatrix}L_{r''} \\ q\end{bmatrix} A\] is of rank $r'' + 1$,
and $S_{r'} = (\Z / p^e \Z)^m$. 
On the other hand, for any $r'  + 1\leq r'' \leq r$ and any $q \in S_{r''} \setminus S_{r'' - 1}$, the inner rank of 
\[\begin{bmatrix}L_{r'' - 1} \\ i\cdot q + x\end{bmatrix}A\]
is $r''$ for any $1 \leq i \leq p - 1$ and $x \in S_{r'' - 1}$. 
%Since $q_{r''} A$ is not in $\rowspan(L_{r'' - 1} A)$ for any $r' \leq r'' \leq r$, $q_{r''} + x \cdot L_{r'' - 1}$ is not in $\rowspan(L_{r'' - 1}A)$ for any non-zero row vector $x$. 
%And thus, $i \cdot q_{r''} +  x \cdot L_{r'' - 1}$ for any $1 \leq i \leq p - 1$ is not in $\rowspan(L_{r'' - 1}A)$ for any non-zero row vector $x$. 
%Thus, we have  
%\[\frac{|S_{r''}|}{|S_{r'' - 1}|} \geq p.\]
%Notice that every $q \in (\Z / p^e \Z)^m \setminus S_{r'}$ satisfies the condition that \[\rank\left(\begin{bmatrix} L \\ q \end{bmatrix}\right) > r'.\] 
Hence, for a randomly sampled row vector $q \in (\Z / p^e \Z)^m$, 
\[\begin{bmatrix}L \\ q\end{bmatrix}A\] is a matrix of inner rank $r' + 1$ with probability at least $1 - (1 / p)^{r - r'}$. 
\end{proof}

\begin{lemma}\label{lem:individualization_single_matrix}
Let $A$ be a matrix in $\M(m\times n,  \Z / p^e \Z)$ with inner rank $r$. 
Given an integer $r' \leq r / 2$ and a parameter $0 < \delta < 1$ such that $\log(1 / \delta) \geq r$, 
let $Q$ be a random matrix in \[\M\left(\left\lceil \frac{8 r' \log (1/\delta)}{r}\right\rceil\times m,  \Z / p^e \Z\right).\] With probability at least $1 - \delta$, $QA$ is of inner rank at least $r'$. 

\end{lemma}
\begin{proof}
%We can write $P$ as a multiplication of $A$ and $B$ for some $A \in M(m, \gamma, \Z / p^e \Z)$ and $B \in M(\gamma, n, \Z / p^e \Z)$ such that the row span of $P$

Let $\alpha = \left\lceil \frac{8 r' \log (1/\delta)}{r}\right \rceil$. 
Without loss of generality,
we assume the matrix $Q$ is sampled sequentially by rows $q_1, . . . , q_\alpha$, where $q_i$ is the $i$-th row of $Q$.
Let $Q_\beta$ denote the matrix $Q$ that contains only the first $\beta$ rows for some $1 \leq \beta \leq \alpha$. 

Suppose that after sampling the first $\beta$ rows for some positive integer $\beta$, $Q_\beta A$ is of inner rank $\gamma$. 
%Then we can write $P$ as a multiplication of $A$ and $B$ for some $A \in M(m, r, \Z / p^e \Z)$ and $B \in M(r, n, \Z / p^e \Z)$ such that there is a matrix $A_{\beta} \in M(\beta, \gamma, \Z / p^e \Z)$ such that $A_\beta B_\gamma$ equals $Q_\beta P$, where $B_\gamma$ equals the first $\gamma$ rows of $B$. 
%Hence for an arbitrary row vector $q \in (\Z / p^e \Z)^n$ such that $q A$ has at least one non-zero entries in coordinates greater than or equal to $\beta$, 
%$\begin{bmatrix}Q_\beta \\ q^t\end{bmatrix}B$ is a matrix of inner rank $\gamma + 1$.
%Since $A$ is full rank, the probability of sampling such a row vector $q$ is at least $1 / p^{r - \gamma}$. 
By Lemma~\ref{lem:sampling_p_to_e}, the probability of sampling a random row vector $q \in (\Z / p^e \Z)^m$ such that \[\begin{bmatrix}Q_\beta \\ q\end{bmatrix}B\] is a matrix of inner rank $\gamma + 1$ is at least $1 - (1 / p)^{r - \gamma}$.
By sampling $\lceil \log (r' / \delta)  / (r - r')\rceil $ new rows of $Q$, the inner rank of $Q$ increases by at
least one with probability
\[1 - \left(\frac{1}{p} \right)^{(r - \gamma) \lceil \log(r' / \delta) / (r - r')\rceil} \geq 1 - \left(\frac{1}{p} \right)^{\log (r' / \delta)}\geq 1 - \frac{\delta}{r'}. \]
On the other hand, observe that 
\[r' \left\lceil \frac{\log (r' / \delta)}{(r - r')}\right\rceil \leq \frac{2 r' \log (r' / \delta)}{r - r'}\leq \frac{8r' \log (1 / \delta)}{r}. \]
By union bound, with probability at least $1 - \delta$, if $Q$ contains at least $\left\lceil \frac{8 r' \log (1/\delta)}{r}\right\rceil$ rows, the inner rank of $QA$ is a matrix of inner rank at least $r'$. 
\end{proof}


\begin{lemma}\label{lem:individualization}
For any $\Z/p^e \Z$-submodule of $\M(m\times n, \Z / p^e \Z)$ containing $\tau$ different matrices and any integer $r \geq 4$, 
denote 
$t = \left\lceil 32 \max\{\log \tau, r \}/ \sqrt{r} \right\rceil$. 
There is a matrix $L \in \M(t\times m, \Z / p^e \Z)$ and a matrix $R \in \M(n\times t, \Z / p^e \Z)$ such that for any matrix $A$ in the submodule of rank at least $r$, $LAR$ is a non-zero matrix. 
\end{lemma}
\begin{proof}
Let $A$ be an arbitrary matrix of rank at least $r$ in the submodule. By Lemma~\ref{lem:individualization_single_matrix} with $r' = \lfloor \sqrt{r}\rfloor$ and $\delta = \min\{1 / 4\tau, 1 / 2^r \}$, for a random matrix $L \in \M(t, m, \Z/p^e \Z)$, $LA$ is of rank at least $r'$ with probability at least $1 - \delta$. 

If this is the case, then by Lemma~\ref{lem:individualization_single_matrix}, with $r' = 1$ and $\delta = \min\{1 / 4\tau, 1 / 2^r \}$, for a random matrix $R \in \M(n, t, \Z/p^e \Z)$, $LAR$ is of rank at least $1$ with probability at least $1 - \delta$. 


By union bound, with constant probability, for random $L$ and $R$, $LAR$ is a non-zero matrix
for all the matrices $A$ in the submodule of rank at least $r$.
\end{proof}



\subsubsection{Low rank matrix submodule characterization over $\Z/p^e \Z$}


We define the rank of a $(\Z / p^e \Z)$-submodule of $(\Z / p^e \Z)^n$ to be the smallest integer $r$ such that 
the submodule has a generating set of size  $r$. %\xnote{Do not know if it is good to call it rank here. Please let me know if there is a more standard term.}]\ynote{This is good due to local ring.}

\begin{lemma}\label{lem:space_dim_prelim}
Let $A_1, A_2, \dots, A_k$ be $k$ matrices in $\M(m\times n, \Z / p^e \Z)$ for some positive integers $k, m, n$. 
If there exist $d$ row vectors
 $x_1, x_2, \dots, x_d \in ( \Z / p^e \Z)^m$ such that for every $1 \leq i < d$, 
% $\left\langle \left\{x_j A_\ell : 1 \leq j \leq i, 1 \leq \ell \leq k\right\}\right\rangle$ is not a subspace of  $\left\langle \left\{x_j A_\ell : 1 \leq j \leq i + 1, 1 \leq \ell \leq k\right\} \right\rangle$.
the following condition holds
\begin{align}\label{equ:rank_bound}
\rank\big(\left\langle \left\{x_j A_\ell : 1 \leq j \leq i, 1 \leq \ell \leq k\right\}\right\rangle \big) \neq \rank\big(\left\langle \left\{x_j A_\ell : 1 \leq j \leq i + 1, 1 \leq \ell \leq k\right\} \right\rangle\big),
\end{align}
then there is a linear combination of $A_1, A_2, \dots, A_k$ with inner rank at least $(1 - 1 / p)d$.
\end{lemma}
\begin{proof}
%We prove that there exist $\alpha_1, \dots, \alpha_k \in \mathbb{F}_p$ such that $\sum_{i=1}^k \alpha_i A_i$ is a matrix of rank at most $()$. 
Let $X$ be the matrix in $\M(d\times n, \Z / p^e \Z)$ such that $x_i$ is the $i$-th row of $X$.
To prove the lemma, it is sufficient to show that if
$\alpha_1, \dots, \alpha_k $ are uniformly and independently  sampled from $\Z / p^e \Z$, then
we have the following expectation estimation. %of the rank of $X (\sum_{\ell=1}^k \alpha_\ell A_\ell)$ is at least $(1 - 1/p)d$.
\[\mathbb{E}_{\alpha_1, \dots, \alpha_k}\left[\rank \left({X \left(\sum_{\ell=1}^k \alpha_\ell A_\ell\right) }\right)\right] \geq \left(1 - \frac{1}{p}\right)d.\]
If this is the case, then there exist $\alpha_1, \dots, \alpha_k \in \Z / p \Z$ such that 
\[\rank{ \left(\sum_{\ell=1}^k \alpha_\ell A_\ell \right)} \geq \rank {\left(X \left(\sum_{\ell=1}^k \alpha_\ell A_\ell\right) \right)} \geq \left(1 - \frac{1}{p}\right)d, \]
and then the lemma follows. 

For any $1 \leq i \leq d$, 
let $D_i$ be a $\Z / p^e \Z$-row span constructed as follows:
Let $D_i$ be \[\left\langle \left\{x_j A_\ell : 1 \leq j \leq i, 1 \leq \ell \leq k\right\}\right\rangle\] initially. 
Find a vector $u = x_i (\sum_{j = 1}^k \alpha_j A_j)$ for some $\alpha_1, \dots, \alpha_k \in \Z / p^e \Z$ such that $u$ is not in $D_i$ but $\rank(\langle \{u\}\cup D_i \rangle) = \rank(D_i)$. 
Let $D_i = \langle \{u\}\cup D_i \rangle$. 
Repeat above process until no vector $u$  can be added to $D_i$ to increase the size of $D_i$. 
Let $S_i$ be the set of row vectors
\[\left\{(\alpha_1, \dots, \alpha_k) \in (\Z / p^e \Z)^k : x_i \left(\sum_{\ell = 1}^k \alpha_\ell A_\ell \right) \in D_{i}\right\}.\]
By Inequality~(\ref{equ:rank_bound}), $S_i$ does not equal to $(\Z / p^e \Z)^k$.


Since $D_{i-1}$ is a $\Z / p^e \Z$-submodule of $(\Z / p^e \Z)^n$, 
%Since $\left\langle \left\{x_j A_\ell : 1 \leq j \leq i - 1, 1 \leq \ell \leq k\right\} \right\rangle$ is a row vector space,
for any $(\alpha_1, \dots, \alpha_k)$ in $S_i$ and $\mu \in \Z / p^e \Z$, $\mu (\alpha_1, \dots, \alpha_k)$ is in $S_i$, and 
if both $(\alpha_1, \dots, \alpha_k)$ and $(\beta_1, \dots, \beta_k)$ are in $S_i$, then $(\alpha_1 + \beta_1, \dots, \alpha_k + \beta_k)$ is also in $S_i$. 
Hence, the row vectors in $S_i$ form a $\Z / p^e \Z$-submodule of $(\Z / p^e \Z)^k$, and thus 
$|S_i| \leq p^{k e-1}$.


On the other hand, for any $(\alpha_1, \dots, \alpha_k) \notin S_i$, 
since $\rank(\langle \{x_i (\sum_{\ell = 1}^k \alpha_\ell A_\ell)\} \cup D_i\rangle) > \rank(D_i)$, 
we have 
\[\rank\left(\begin{bmatrix} x_1 \\ \vdots \\ x_i \end{bmatrix}\left(\sum_{\ell=1}^k \alpha_\ell A_\ell\right)\right) > \rank\left(\begin{bmatrix} x_1 \\ \vdots \\ x_{i-1} \end{bmatrix}\left(\sum_{\ell=1}^k \alpha_\ell A_\ell\right)\right).\]
Thus,
\[\mathrm{Pr}_{\alpha_1, \dots, \alpha_k}\left[\rank\left(\begin{bmatrix} x_1 \\ \vdots \\ x_i \end{bmatrix}\left(\sum_{\ell=1}^k \alpha_\ell A_\ell\right)\right) > \rank\left(\begin{bmatrix} x_1 \\ \vdots \\ x_{i-1} \end{bmatrix}\left(\sum_{\ell=1}^k \alpha_\ell A_\ell\right)\right)\right] \geq \frac{p^{ek} - |S_i|}{p^{ek}} \geq 1 - \frac{1}{p}.\]

Thus, we have \begin{align*}& \mathbb{E}_{\alpha_1, \dots, \alpha_k}\left[\rank {X \left(\sum_{\ell=1}^k \alpha_\ell A_\ell\right) }\right]  \\ \geq & 
\sum_{i=1}^d \mathrm{Pr}_{\alpha_1, \dots, \alpha_k}\left[\rank\left(\begin{bmatrix} x_1 \\ \vdots \\ x_i \end{bmatrix}\left(\sum_{\ell=1}^k \alpha_\ell A_\ell\right)\right) > \rank\left(\begin{bmatrix} x_1 \\ \vdots \\ x_{i-1} \end{bmatrix}\left(\sum_{\ell=1}^k \alpha_\ell A_\ell\right)\right)\right]  \\
\geq & \left(1 - \frac{1}{p}\right)d.\qedhere\end{align*}



\begin{comment}
For any $1 \leq i \leq d$, 
let $D_i$ be $\left\langle \left\{x_j A_\ell : 1 \leq j \leq i, 1 \leq \ell \leq k\right\}\right\rangle$,  
%a submodule of $(\Z / p^e \Z)^n$ generated by $rank(\left\langle \left\{x_j^t A_\ell : 1 \leq j \leq i, 1 \leq \ell \leq k\right\}\right\rangle )$ elements such that $D_i$ contains both $\left\langle \left\{x_j^t A_\ell : 1 \leq j \leq i, 1 \leq \ell \leq k\right\}\right\rangle$ and $D_{i-1}$, 
and $S_i$ be the set of row vectors
\[\left\{(\alpha_1, \dots, \alpha_k) \in (\Z / p^e \Z)^k : x_i \left(\sum_{\ell = 1}^k \alpha_\ell A_\ell \right) \in D_{i-1}\right\}.\]
By Inequality~(\ref{equ:rank_bound}), $S_i$ does not equal to $(\Z / p^e \Z)^k$.

Since $D_{i-1}$ is a $\Z / p^e \Z$-submodule of $(\Z / p^e \Z)^n$, 
%Since $\left\langle \left\{x_j A_\ell : 1 \leq j \leq i - 1, 1 \leq \ell \leq k\right\} \right\rangle$ is a row vector space,
for any $(\alpha_1, \dots, \alpha_k)$ in $S_i$ and $\mu \in \Z / p^e \Z$, $\mu (\alpha_1, \dots, \alpha_k)$ is in $S_i$, and 
if both $(\alpha_1, \dots, \alpha_k)$ and $(\beta_1, \dots, \beta_k)$ are in $S_i$, then $(\alpha_1 + \beta_1, \dots, \alpha_k + \beta_k)$ is also in $S_i$. 
Hence, the row vectors in $S_i$ form a $\Z / p^e \Z$-submodule of $(\Z / p^e \Z)^k$.
%By Inequality~(\ref{equ:rank_bound}), $S_i$ does not contain all the vectors in $\mathbb{F}_p^k$.
Hence $|S_i| \leq p^{k e-1}$.
We have for each $1 \leq i \leq d$, 
\[\mathrm{Pr}_{\alpha_1, \dots, \alpha_k}\left[x_i \left(\sum_{\ell=1}^k \alpha_\ell A_\ell\right) \notin  D_{i-1}\right] = \frac{p^{ek} - |S_i|}{p^{ek}} \geq 1 - \frac{1}{p}.\]
If $x_i \left(\sum_{\ell=1}^k \alpha_\ell A_\ell\right)$ is not in $D_{i-1}$, 
then 
\[\rank\left(\begin{bmatrix} x_1 \\ \vdots \\ x_i \end{bmatrix}\left(\sum_{\ell=1}^k \alpha_\ell A_\ell\right)\right) > \rank\left(\begin{bmatrix} x_1 \\ \vdots \\ x_{i-1} \end{bmatrix}\left(\sum_{\ell=1}^k \alpha_\ell A_\ell\right)\right)\]
% $x_i \left(\sum_{\ell=1}^k \alpha_\ell A_\ell\right)$ is not a linear combination of $x_1 \left(\sum_{\ell=1}^k \alpha_\ell A_\ell\right), \dots,  x_{i-1} \left(\sum_{\ell=1}^k \alpha_\ell A_\ell\right)$. 
Thus, we have \begin{align*}& \mathbb{E}_{\alpha_1, \dots, \alpha_k}\left[\rank {X \left(\sum_{\ell=1}^k \alpha_\ell A_\ell\right) }\right]  \\ \geq & 
\sum_{i=1}^d \mathrm{Pr}_{\alpha_1, \dots, \alpha_k}\left[x_i \left(\sum_{\ell=1}^k \alpha_\ell A_\ell\right) \notin D_{i-1}\right]  \\
\geq & \left(1 - \frac{1}{p}\right)d.\end{align*}
\end{comment}
\end{proof}

\begin{lemma}\label{lem:char_main}
Let $\mathfrak{A}$ be a $\Z / p^e \Z$-submodule of $M(m, n, \Z / p^e \Z)$ 
%or a skew-symmetric matrix subspace of $M(n, n, \Z / p^e \Z)$
such that for each matrix $A$ in the submodule, the inner rank of $A$ is at most $r$ for some positive integer $r$. 
There exists a matrix $P \in \GL(m, \Z / p^e \Z)$ and an integer $d = O(er)$ such that 
\[\rank\left(\left\langle \bigcup_{i = d+1}^{m} \langle p_i \mathfrak{A}\rangle \right\rangle\right) \leq O(e^2 r^2),\]
where $p_i$ denotes the $i$-th row of $P$.

%There exists a set of matrices $O(e^2 r^2)$ row vectors $S$ such that 
%$\{x \in (\Z / p\Z)^m : \langle x A : A \in \mathfrak{A}\rangle \subseteq \langle S \rangle\}$ is a subspace of $\Z / p^e \Z$ of dimension at least $n - O(e^2 r^2)$ and 
%there exists a set of $O(e^2 r^2)$ row vectors $S'$ such that 
%$\{x \in (\Z / p\Z)^m : \langle x A : A \in \mathfrak{A}\rangle \subseteq \langle S \rangle\} \cup S'$ generates $\Z / p^e \Z$.

%There exists two matrices $L \in \GL(m, \Z / p^e \Z)$ and $R \in \GL(n, \Z / p^e \Z)$ such that for every $A \in \mathfrak{A}$, $LAR$ has non-zero entries only in the first  $O(??r^2)$ rows and columns. 
%set $\Lambda$ of linear independent row vectors in $\mathbb{F}_p^n$ with $|\Lambda| = O(r^2)$ such that 
%\begin{enumerate}
%\item 
%$\dim{\ker(\mathfrak{A}, \Lambda)} \geq m - O(??r)$ if $\mathfrak{A}$ is a matrix subspace of $M(m, n, \Z / p^e \Z)$, or
%\item 
%$\dim{\ker_\mathrm{skew}(\mathfrak{A}, \Lambda)} \geq n - O(??r^2)$ if $\mathfrak{A}$ is a skew-symmetric matrix subspace of $\SS(n, \Z / p^e\Z)$.
%\end{enumerate}
\end{lemma}

\begin{proof}
%We first consider the case that $\mathfrak{A}$ is a matrix subspace of $M(m, n, \Z / p^e \Z)$.
Let $A_1, \dots, A_k$ be a generating set of $\mathfrak{A}$.
Let $x_1, x_2, \dots, x_d \in (\Z / p^e \Z )^m$ be $d$ row vectors for some positive integer $d$ 
 such that there is a matrix in $\GL(m, \Z / p^e \Z)$ with $x_i$ being $i$-th row, %$x_i$ is not a multiple of $p$ to another vector in $(\Z / p \Z )^m$, 
 and for any $1 \leq i \leq d$, there are at least $p^{ek} / 2$ different $\alpha_1, \alpha_2, \dots, \alpha_k \in \Z / p^e \Z$ satisfying
\begin{equation}\label{equ:characterization}\rowspan\left(X_{i-1} \left(\sum_{j =1}^k \alpha_j A_j\right)\right) \subsetneq \rowspan\left(X_i \left(\sum_{j =1}^k \alpha_j A_j\right)\right)\end{equation} where for each $1 \leq i \leq d$,
$X_i$ is the $i \times m$ matrix with $x_j$ as the $j$-th row of $X_i$ for all $1 \leq j \leq i$.
Since every linear combination $\sum_{j =1}^k \alpha_j A_j$ of $A_1, \dots, A_k$ has at most $e \cdot r$ different $i$ such that (\ref{equ:characterization}) holds, 
we have 
\[d \cdot \frac{p^{ek}}{2} \leq e \cdot r \cdot p^{ek},\]
which implies $d \leq 2er$. Suppose there does not exist a row vector $x_{d+1} \in (\Z / p^e \Z)^m$ such that there are at least $p^{ek} / 2$ different $\alpha_1, \alpha_2, \dots, \alpha_k\in\Z / p^e \Z$ satisfying
\[\rowspan\left(X_{d} \left(\sum_{j =1}^k \alpha_j A_j\right)\right) \subsetneq \rowspan\left(X_{d+1} \left(\sum_{j =1}^k \alpha_j A_j\right)\right).\]
%\[\rank{X_{d+1} \left(\sum_{j =1}^k \alpha_j A_j\right)} > \rank{X_d \left(\sum_{j =1}^k \alpha_j A_j\right)}.\] 
Then for each row vector $v \in (\Z / p\Z)^m$ such that $v\notin \langle x_1, \dots, x_d\rangle$, 
there exist $\beta_1, \dots, \beta_d \in \Z / p^e\Z$ such that the following condition holds for at least $\frac{p^{ek}}{ 2p^{ed}}$ different $\alpha_1, \alpha_2, \dots, \alpha_k \in \Z / p^e\Z$
\[v \cdot \left(\sum_{i = 1}^k \alpha_k A_k\right) = \left(\sum_{i=1}^d \beta_i x_i\right)\left(\sum_{i = 1}^k \alpha_k A_k\right).\]
Thus, the rank of $\langle (v -  \sum_{i=1}^d \beta_i x_i) \mathfrak{A}\rangle$ is at most $ed + 1 = O(e^2r)$. 
By Lemma~\ref{lem:space_dim_prelim}, one can construct a matrix $P \in \mathrm{GL}(m, \Z / p^e \Z)$ satisfying the following properties: Let $p_i$ denote the row vector of the $i$-th row of $P$.
\begin{enumerate}
\item $p_i = x_i$  for any $1 \leq i \leq d$.
\item $\langle \cup_{i = d + 1}^m \cup_{j = 1}^k\{ p_i A_j \}\rangle$ is of rank at most $O(e^2r^2)$. 
\end{enumerate} 
%Hence, for given $x_1, \dots, x_d$, there is an invertible $n \times n$ matrix $X$ such that for any $1 \leq i \leq d$, $i$-th row of $X$ is $x_i$ and for any $d + 1 \leq i \leq n$, $<X_i A_j>_{j=1}^k$ spans a space of dimension at most $O(t)$. 
Any such matrix $P$ proves the lemma.
%\ynote{Proof incomplete here?}
\end{proof}

\begin{lemma}\label{lem:char_skew}
Let $r$ be a positive integer and $\mathfrak{A}$ be a $\Z / p^e \Z$-submodule of $\M(n, \Z / p^e \Z)$ 
%or a skew-symmetric matrix subspace of $M(n, n, \Z / p^e \Z)$
such that 
for each matrix $A$ in the submodule, $A$ is an alternating matrix and the inner rank of $A$ is at most $r$. % for some positive integer $r$. 
Then there exists an integer $t = O(e^2 r^2)$ and a matrix $P \in \GL(n, \Z / p^e \Z)$ such that for each matrix $A \in \mathfrak{A}$,
$i$-th row $j$-th column of $P A P^t$ is zero for any $i \geq t, j \geq t$. % for some $t = O(e^2 r^2)$. 
\end{lemma}
\begin{proof}
By Lemma~\ref{lem:char_main}, there exists an integer $d = O(er)$ and a matrix $P_0 \in \GL(n, \Z / p^e \Z)$ such that  
\[\rank\left(\left\langle \bigcup_{i = d+1}^{m} \langle p_{0, i} \mathfrak{A}\rangle \right\rangle\right) \leq O(e^2 r^2),\]
where $p_{0, i}$ denotes the $i$-th row of $P_0$. 
Let $S$ be a smallest generating set of  $\left\langle \bigcup_{i = d+1}^{m} \langle p_{0, i} \mathfrak{A}\rangle \right\rangle$. We have $|S| = O(e^2 r^2)$. 
Denote $S = \{s_1, \dots, s_{|S|}\}$. 

Let $t_0 = d$. 
We perform the following process for $i = 1, \dots, |S|$: 
\begin{enumerate}
\item Find a  $j \in \{t_{i-1} + 1, \dots n\}$ such that $p_{i-1, j} \cdot s_i^t \in \Z / p^e \Z$ (suppose $p_{i - 1, j}$ denote $j$-th row of $P_{i-1}$) is non-zero and generates all the $p_{i-1, k} \cdot s_i^t$ under additions for any $t_{i-1} + 1\leq k \leq n$. %~\xnote{Here, I suppose every subring of $\Z / p^e \Z$ contains a single generator, correct me if I am wrong.}. 
If there is no such $j$, then let $P_i = P_{i-1}$, $t_i = t_{i - 1}$,  and continue to the process for $i + 1$. 
\item Otherwise, let $t_i = t_{i-1} + 1$ and $P_i'$ be the resulted matrix by switch $j$-th row and $t_i$-th row of $P_{i-1}$.
\item Otherwise, $P_i$ be the matrix such that the first $t_i$ rows are the same as $P_i'$, and for any $t_i + 1 \leq j \leq n$, 
the $j$-th row of $P_i$ is a linear combination of $j$-th row and $t_i$-th row of $P_i'$ such that its inter product with $s_i$ is zero. 
\end{enumerate}
Let $P$ be $P_{|S|}$. 
It is easy to verify that $t_{|S|} = O(e^2 r^2)$, and for each matrix $A \in \mathfrak{A}$,
$i$-th row $j$-th column of $P A P^t$ is zero for any $i \geq t_{|S|}, j \geq t_{|S|}$.
In addition, since $P$ is obtained by applying multiple row switching operations and row addition operations to $P_0$, $P$ is also in $\GL(n, \Z / p^e \Z)$. 
\end{proof}

\subsubsection{Reduction to matrix tuple isometry}
\begin{comment}
Let $T$ be an $n \times n \times m$ tensor over $\Z / p^{e_0} Z$. Let $\mathfrak{X}_{T, i}$ be the $n \times m$ matrix over $\Z / p^{e_0} \Z$ such that $\mathfrak{X}_{T, i}[k, j] = T[i, j, k]$ for any $1 \leq j \leq n, 1 \leq k \leq m$, and $\mathfrak{Z}_{T, k}$ be the $n \times n$ matrix over $\Z / p^{e_0} \Z$ such that $\mathfrak{Z}_{T, k}[i, j] = T[i, j, k]$ for any $1 \leq i, j \leq n$.

Let $\alpha_1, \dots, \alpha_{e_0}$ be the integers such that the abelian group generated by $\mathfrak{Z}_{T, k}$ for all the $k$ is isomorphic to 
\[\bigoplus_{\ell = 1}^{e_0} \left(\Z / p^{\ell } \Z\right)^{\alpha_\ell},\]
and $\beta_1, \dots, \beta_{e_0}$ be the integers such that the abelian group generated by $\mathfrak{X}_{T, i}$ for all the $i$ is isomorphic to 
\[\bigoplus_{\ell = 1}^{e_0} \left(\Z / p^{\ell } \Z\right)^{\beta_\ell}.\]
\end{comment}
We define some additional notations used in this section. 
Let $\tA$ be an $n \times m \times h$ tensor over $\Z / p^e \Z$. 
Let $X = (x_{i', i})\in \M(n'\times n, \Z/ p^e \Z), Y = (y_{j', j}) \in \M(m'\times m, \Z/ p^e \Z)$, and $Z = (z_{k', k}) \in \M(h'\times h, \Z / p^e \Z)$. 
We define $(X, Y, Z)\circ \tA$ to be the $n' \times m' \times h'$ dimensional tensor $\tB = (b_{i', j' ,k'})$ such that
$b_{i', j', k'} = \sum_{i, j, k} x_{i', i} y_{j', j} z_{k', k} a_{i, j, k}$. 
Here, $X$, $Y$, and $Z$ do not need to be square matrices.  



For a row vector $x = (x_1, \dots, x_n) \in (\Z / p^{e} \Z)^n$, we denote $\mathfrak{X}_{\tA, x} = \sum_{\ell = 1}^n x_\ell \cdot \mathfrak{X}_{\tA, \ell}$.
For a row vector $y = (y_1, \dots, y_m) \in (\Z / p^{e} \Z)^m$, we denote $\mathfrak{Y}_{\tA, y} = \sum_{\ell = 1}^m y_\ell \cdot \mathfrak{Y}_{\tA, \ell}$.
For a vector $z = (z_1, \dots, z_h) \in (\Z / p^{e} \Z)^h$, we denote $\mathfrak{Z}_{tA, z} = \sum_{\ell = 1}^h z_\ell \mathfrak{Z}_{\tA, \ell}$.
The order of a matrix $\mathfrak{X}_{\tA, x}$, denoted as $\ord(\mathfrak{X}_{\tA, x})$, is the smallest non-negative integer $t$ such that 
$t \cdot \mathfrak{X}_{\tA, x}$ is a zero matrix. 
%We say a matrix $\mathfrak{X}_{\tA, x}$ is of order $t$ if $\mathfrak{X}_{\vA, x}$ is an element of order $t$ in the abelian group generated by $\mathfrak{X}_{\tA, i}$ for all the $i$. 
The orders of matrices $\mathfrak{Y}_{\tA, y}$ and $\mathfrak{Z}_{\tA, z}$ are defined in the same way. 

Let $O_{\mathfrak{X}, t} = \{x \in (\Z / p^{e} \Z)^n : \ord(\mathfrak{X}_{T, x}) \leq t\}$. Clearly, $\langle O_{\mathfrak{X}, t} \rangle = O_{\mathfrak{X}, t}$ for any $t$.
We also define $O_{\mathfrak{Y}, t}$ and $O_{\mathfrak{Z}, t} $  in the same way.

\begin{lemma}\label{lem:enumeration}
Let $\tA$ be an alternating $n \times n \times m$ tensor over $\Z / p^{e} \Z$, and $e_0$ be an integer such that $1 \leq e_0 \leq e$. %Let $t$ be the rank of $O_{\mathfrak{X}, \gamma} $. %, and $X_e$ be an arbitrary $M(t, n, \Z / p^{e_0} \Z)$ such that $rowspan(X_e) = O_{\mathfrak{X}, p^e}$.
%\ynote{Added:} 
Suppose the abelian group generated by $\mathfrak{X}_{\tA, x}$ for all the $x \in (\Z / p^{e}\Z)^n$ is isomorphic to 
\[\bigoplus_{\ell = 1}^{e} \left(\Z / p^{\ell } \Z\right)^{\beta_\ell}.\]
For any integer $4 \leq r \leq n$, 
there are matrices $X_0, X_1, X_2, X_3$ and $T$ satisfying the following conditions:
\begin{enumerate}
%\item For every row vector $s \in S$, every entry of $p^{e_0} \cdot s$ is zero, and 
%$|S| = O(e_0^2 r^2 + \sum_{i = e_0 + 1}^{e} \beta_i)$.
\item Every row of $X_0, X_1, X_2, X_3$ and $T$ is in $(\Z / p^{e} \Z)^n$. In addition, 
\[\rowspan\left(\begin{bmatrix} X_0\\ X_1 \\ X_2 \\ X_3\end{bmatrix}\right) = (\Z / p^{e} \Z)^n,\]
and 
\[\rowspan\left(\begin{bmatrix} X_1 \\ X_2 \\ X_3\end{bmatrix}\right) = O_{\mathfrak{X}, p^{e_0}}.\]
\item $T$ has $O(e_0^2 r^2)$ rows from $O_{\mathfrak{X}, p^{e_0}}$. % such that 



\item $X_0$ has $\sum_{i = e_0+ 1}^{e} \beta_i$ rows, $X_1$ has 
$O(\max\{\log(p^{\sum_{i = 1}^{e} i \cdot \alpha_i}), r\} /  \sqrt{r} )$ rows, and
$X_2$ has $O(e_0^2 r^2 + \sum_{i = e_0 + 1}^{e} \beta_i)$ rows. 
Rows of $X_3$ correspond to a smallest generating set of 
\[\left\langle \left\{x \in O_{\mathfrak{X}, p^e} : \forall s \in S, s \cdot x^t = 0 \text{ and } 
x \cdot \mathfrak{Z}_{\tA, z} \in S \text{ for all } z \in (\Z / p^e \Z)^m \text{ s.t. } X_1  \mathfrak{Z}_{\tA, z} X_1^t = 0,   \right\}\right\rangle,\]
where \[S = \langle \{s  \in (\Z / p^e \Z)^m : \exists z  \text{ s.t. }X_1 \mathfrak{Z}_{\tA, z} X_1^t = 0, s \in \rowspan(T \mathfrak{Z}_{\tA, z})\}\rangle.\]
%$S$ has a generating set of size at most $O(e_0^2 r^2 + )$.

\item For any $z \in (\Z / p^e \Z)^m$ such that 
\[\begin{bmatrix} X_0 \\ X_1 \\ X_2 \end{bmatrix} \mathfrak{Z}_{\tA, z} \begin{bmatrix} X_0 \\ X_1  \\ X_2 \end{bmatrix}^t= 0,\]
$X_3 \mathfrak{Z}_{\tA, z}  X_3^t = 0$. 

%\item $U_0$ contains $O(\lceil 32 \max\{d \cdot \log(p^e), r\}/ \sqrt{r} \rceil)$ rows, $U_1$ contains $O(e^2 r^2)$ rows, and $U_2$ contains at most $n$ rows.
\end{enumerate} 
\end{lemma}
\begin{proof}
Since the abelian group generated by $\mathfrak{X}_{\tA, x}$ over $x \in (\Z / p^{e}\Z)^n$ is isomorphic to $\bigoplus_{\ell = 1}^{e} \left(\Z / p^{\ell } \Z\right)^{\beta_\ell},$
there is an invertible matrix $X \in \GL(n, \Z / p^{e}\Z)$ such that for any $1 \leq i \leq e$ and any $j\in[\sum_{k = i + 1}^{e} \beta_k + 1, \sum_{k = i}^{e} \beta_k]$,
%\ynote{Was: $(\sum_{k = i + 1}^{e} \beta_k) + 1 \leq  j  \leq \sum_{k = i}^{e} \beta_k$},
$\mathfrak{X}_{\tA, j}$ is of order $p^i$, and for $j > \sum_{\ell = 1}^{e} \beta_\ell$,   $\mathfrak{X}_{\tA, j}$ is the zero matrix.
Let $X_0$ be the first $\sum_{\ell = e_0 + 1}^{e} \beta_\ell$ rows of $X$. 

Let $X'$ be a matrix such that the rows of $X'$ correspond to a smallest generating set of $O_{\mathfrak{X}, p^{e_0}}$, and $n'$ denote the number of rows in $X'$. We have $n' \leq n$ and $\langle \rowspan(U_0) \cup \rowspan(X')\rangle = (\Z / p^e \Z)^n$.
For any $x \in \rowspan(X')$, $\mathfrak{X}_{\tA, x}$ is a matrix of order at most $p^{e_0}$. 
%Let  $X'$ be the matrix obtained as follows: For any $e_0 + 1 \leq i \leq e$ and any $(\sum_{k = i + 1}^{e_0} \beta_k) + 1 \leq  j  \leq \sum_{k = i}^{e_0} \beta_k$, $j$-th row of $X'$ equals $j$-th row of $X$ multiply $p^{i - e}$. 
%For $j > \sum_{k = e + 1}^{e_0} \beta_k$, the $j$-th row of $X'$ equals $j$-th row of $X$.
%By the construction of $X'$, the row span of $X'$ equals $O_{\mathfrak{X}, p^e}$. 
%Hence, 
%Since  the row span of $X'$ equals $O_{\mathfrak{X}, p^e}$, 
$(X', X', I) \circ \tA$ has every entry be a multiple of $p^{e - e_0}$. 


Let $\tA'$ be the tensor over $\Z / p^{e_0} \Z$ equivalent to $(X', X', I) \circ \tA$. 
By Lemma~\ref{lem:individualization}, %there exists a set of $O(\max\{d \cdot \log(p^e), r\}/ \sqrt{r} \rceil)$ row vectors such that 
there is a matrix 
\[U_{1}' \in \M\left(O\left(\frac{\max\left\{\log(p^{\sum_{i = 1}^{e} i \cdot \alpha_i}), r\right\}}{ \sqrt{r}}\right), n', p^{e_0}\right)\] such that 
%every entry of $U_{1}$ is in $\{0, 1, \dots, p^{e_0} - 1\}$, and 
for any matrix $\mathfrak{Z}_{\tA', z}$ of inner rank at most $r$, $U_{1}' \mathfrak{Z}_{\tA', z} U_{1}'^t$ is a non-zero matrix. 

Let $\mathfrak{A}$ be the $\Z / p^{e_0}\Z$-submodule of $\M(n', \Z / p^{e_0} \Z)$ spanned by all the $\mathfrak{Z}_{\tA', z}$ such that $U_{1}'\mathfrak{Z}_{\tA’, z} U_{1}'^t = 0$. 
By Lemma~\ref{lem:char_skew}, there exists a matrix $P \in \GL(n, \Z / p^{e_0} \Z)$ such that the $i$-th row $j$-th column of $P A P^t$ is zero for all the $A \in \mathfrak{A}$ and $i, j \geq O(e_0^2r^2)$. 

Hence, let $P'$ be the matrix of $P$ removing the first $O(e_0^2 r^2)$ rows. 
%By treating both $P$ and $P'$ as matrices in 
We have \[\left\langle \left\{ s \in (\Z / p^{e_0}\Z)^{n'} : \exists z \in (\Z / p^{e_0} \Z)^m \text{ s.t. } U_{1}' \mathfrak{Z}_{\tA', z} U_{1}'^t = 0 \text{ and } s \in \rowspan(P'  \mathfrak{Z}_{\tA', z}) \right\} \right\rangle\] has a generating set of $O(e_0^2 r^2)$ row vectors such that every element in the generating set multiple $p^e$ equals the zero vector. 


Let $U_{1}$ be the matrix over $\Z / p^e \Z$ which has the same dimension as $U_{1}'$ such that every entry of $U_{1}$ equals $U_{1}'$.
By the correspondance between $(X', X', I) \circ \tA$ and $\tA'$,  
there exist matrices $T_0$,
$U_{2}$ and $U_3$ over $\Z / p^e \Z$ all with $n'$ columns satisfying the following conditions:
\begin{enumerate}
\item \[\rowspan\left(\begin{bmatrix} U_{1} \\ U_{2} \\ U_3 \end{bmatrix}\right) = (\Z /  p^e \Z)^{n'}\]
\item $T_0$ contains $O(e_0^2 r^2)$ rows, $U_1$ contains $O(\max\{\log(p^{\sum_{i = 1}^{e} i \cdot \alpha_i}), r\} / \sqrt{r})$ rows, and $U_2$ contains $O(e_0^2 r^2)$ rows. 
\item 
Let 
\[\begin{split}S_0 = \langle \{s \in (\Z / p^{e}\Z)^{n'} : \exists  z \in (\Z / p^{e}\Z)^m \text{ s.t. }& U_1 \mathfrak{Z}_{(X', X', I) \circ \tA, z} U_1^t = 0 \text{ and } \\ & s \in \rowspan(T_0 \mathfrak{Z}_{(X', X', I) \circ \tA, z})\}\rangle.\end{split}\]
%\rowspan(T_0 \mathfrak{Z}_{(X', X', I) \circ \tA, z}) : \text{ for all } z  \text{ s.t. }U_1 \mathfrak{Z}_{(X', X', I) \circ \tA, z} U_1^t = 0\}\rangle$. 
$S_0$ has a generating set of size at most $O(e_0^2 r^2)$.
\item 
The rows of $U_3$ is a generating set of 
\[\begin{split} \langle \{u \in (\Z / p^e \Z)^{n'} : %\forall s \in S_0, s \cdot u^t = 0 \text{ and }  
u \cdot \mathfrak{Z}_{(X', X', I) \circ \tA, z} \in S_0  \text{ for all } 
z \in (\Z / p^e \Z)^m \text{ s.t. }  U_1 \mathfrak{Z}_{(X', X', I) \circ \tA, z} U_1^t = 0 \}\rangle.\end{split}\]
\end{enumerate}


Now we construct $T, X_1, X_2, X_3$ as required by the lemma. 
Let %$T = T_0 X'$, 
$X_1 = U_1 X'$, 
\[S =  \langle \{ s \in (\Z / p^e \Z)^n : \exists z \in (\Z / p^e \Z)^m \text{ s.t. } X_1 \mathfrak{Z}_{\tA, z} X_1^t = 0 \text{ and } s \in \rowspan(T_0 X' \mathfrak{Z}_{\tA, z})\} \rangle. \]
%\rowspan(T \mathfrak{Z}_{\tA, z}) : \text{ for all } z  \text{ s.t. }X_1 \mathfrak{Z}_{\tA, z} X_1^t = 0\}\rangle.\]
By the relation between $\tA$ and $(X', X', I) \circ \tA$, 
$S_1$ has a generating set of size at most $O(e_0^2 r^2 + \sum_{\ell = e_0 + 1}^e \beta_{\ell})$, and there is a matrix $T$ of $O(e_0^2 r^2 + \sum_{\ell = e_0 + 1}^e \beta_{\ell})$ rows such that 
\[S = \langle \cup_{z \in (\Z / p\Z)^m : X_1 \mathfrak{Z}_{\tA, z}X_1^t = 0}\{\rowspan(T \mathfrak{Z}_{\tA, z})\}\rangle.\]

So let $X_3$ be a smallest generating set of 
\[\langle \{x \in O_{\mathfrak{X}, p^{e_0}} : \forall s \in S, s \cdot x^t = 0 \text{ and } x \cdot \mathfrak{Z}_{ \tA, z} \in S  \text{ for all } 
z \in (\Z / p^e \Z)^m \text{ s.t. }  X_1 \mathfrak{Z}_{\tA, z} X_1^t = 0\}\rangle,\]
and $X_2$ contains all the rows of $U_2 X'$ such that $\rowspan(U_3X')$ is a subset of $\langle \rowspan(X_2) \cup \rowspan(X_3)\rangle$.
There exists such a $X_2$ with $O(e_0^2r^2 + \sum_{\ell = e_0 + 1}^e \beta_\ell)$ rows. 
 

%\[\langle \{\rowspan(P' X' \mathfrak{Z}_{T, z}) : \text{for } z \text{ s.t. } U_{1}X' \mathfrak{Z}_{\tA, z} X'^t U_{1}^t = 0 \} \rangle\] has a generating set of $O(e^2 r^2 + \sum_{i = e+1}^{e_0} \beta_i)$ row vectors such that every element in the generating set multiple $p^e$ equals zero vector. 
%Let $S$ be a set of row vectors of size $O(e^2 r^2 + \sum_{i = e+1}^{e_0} \beta_i)$ such that $\langle S \rangle = \langle \{\rowspan(P' X' \mathfrak{Z}_{T, z}) : \text{for } z \text{ s.t. } U_{1, 1}X' \mathfrak{Z}_{T, z} X'^t U_{1, 1}^t = 0 \} \rangle$.

The $X_0, X_1, X_2, X_3$ and $T$ defined above satisfies the requirement of the lemma. 
\begin{comment}

Let $U_{1, 2}$ be the first $\sum_{i = e + 1}^{e_0} \beta_i$ rows of $X$. 
Let $U_{1, 3}$ be the first $O(r^2 e^2)$ rows of $PX'$. 
Let $U_{1, 4}$ be a set of at most $O(e^2 r^2 + \sum_{i = e+1}^{e_0} \beta_i)$ row vectors such that $\langle \rowspan(U_{1, 4}) \cup F \rangle$ contains the row span of $P' X'$, where \[F = \langle \{x \in \rowspan(X') : \forall s \in S, s\cdot x^t = 0 \text{ and } \forall z \text{ s.t. } U_{1, 1} X' \mathfrak{Z}_{T, z} X'^t U_{1, 1}^t = 0,  x \mathfrak{Z}_{T, z} \in \langle S\rangle \} \rangle.\]


Let $U_0 = \begin{bmatrix} U_{1, 1} X' \\ U_{1, 2} \\ U_{1, 3} \\ U_{1, 4} \end{bmatrix}$ and $U_{2}$ be the matrix such that the rows of $U_1$ correspond to a smallest generating set of 
\[F' = \langle  \{x \in (\Z / p^{e_0} \Z)^n : 
\mathfrak{X}_{T, x}\cdot p^e = 0, 
\forall s \in S, s\cdot x^t = 0 \text{ and } \forall z \text{ s.t. } \begin{bmatrix} U_{0} \\ U_1 \end{bmatrix}\mathfrak{Z}_{T, z} \begin{bmatrix} U_0 \\ U_1 \end{bmatrix}^t = 0,  x \mathfrak{Z}_{T, z} \in \langle S\rangle \}\rangle.\]

Clearly, $F'$ contains $F$. Hence, the row span of $\begin{bmatrix}U_0 \\ U_1 \\  U_2 \end{bmatrix}$ is $(\Z / p^{e_0} \Z)^n$. 
And for every two row vectors $u, u' \in U_1$ and $z$ such that $\begin{bmatrix} U_{0} \\ U_1 \end{bmatrix} \mathfrak{Z}_{T, z} \begin{bmatrix} U_{0} \\ U_1 \end{bmatrix}^t = 0$, 
$u \mathfrak{Z}_{T, z} u'^t = s \cdot u'^t = 0$ for some $s \in \langle S \rangle$. 
\end{comment}
\end{proof}

\begin{lemma}\label{lem:lexical_order}
Let $G$ be an abelian $p$-group, and let $\rho$ be a complete strict order of all the elements in $G$. 
Suppose comparing the order of two elements under $\rho$ takes constant time, and 
$G$ is isomorphic to 
$\bigoplus_{\ell = 1}^{e} \left(\Z / p^{\ell } \Z\right)^{k_\ell}$.
There is an algorithm with running time polynomial in the order of $G$ to compute $g_1, \dots, g_{\sum_{\ell = 1}^e k_\ell}$ satisfying the following conditions:
\begin{enumerate}
	\item For any $1 \leq \ell \leq e$, $g_{(\sum_{j = 1}^{\ell - 1} k_j) + 1}, \dots, g_{\sum_{j = 1}^{\ell} k_j}$ are of order $p^\ell$. 
	\item $g_1, \dots, g_{\sum_{\ell = 1}^e k_\ell}$ form a generating set of $G$. 
	\item $g_1, \dots, g_{\sum_{\ell = 1}^e k_\ell}$ has the smallest order, i.e., for any $g_1', \dots, g_{\sum_{\ell = 1}^e k_\ell}'$ satisfying the first two conditions but different to $g_1, \dots, g_{\sum_{\ell = 1}^e k_\ell}$, there is an $1 \leq i \leq \sum_{\ell = 1}^e k_\ell$ such that $g_j' = g_j$ for all the $1 \leq j < i$ and $\tau(g_i) < \tau(g_i')$.
\end{enumerate}
\end{lemma}
\begin{proof}
Here is the algorithm:
\begin{enumerate}
\item Let $\ell = 1$, $i = 0$.
\item Repeat until $g_1, \dots, g_i$ generate $G$:
\begin{enumerate}
	\item Find the element $g_i$ satisfying the following conditions:
	\begin{enumerate}
		\item $g_i$ is of order $p^\ell$, not generated by $g_1, \dots, g_{i-1}$, and there does not exist $g' \in G$ such that $g'^t = g_i$ for some positive integer $t$.
		\item For any other $g_i'$ satisfying the first condition, $\rho(g_i) < \rho(g_i')$. 
	\end{enumerate}
	\item If $g_i$ exists, then $i = i + 1$, otherwise $\ell = \ell + 1$. 
\end{enumerate} 
\end{enumerate}
The correctness of the algorithm is obtained by the characterization of the automorphism of abelian groups.
The running time of the algorithm is obtained by the fact that every group has a generating set of size at most logarithmic in the group order.
\end{proof}

\begin{definition}\label{def:semi_canonical}
Given an $n\times n \times m$ alternating tensor $\tA$ over $\Z / p^e \Z$ and matrices $T, X_0, X_1, X_2$ with rows in $(\Z / p^e \Z)^n$ satisfying Lemma~\ref{lem:enumeration}, 
a tensor $\tB$ is a semi-canonical form of $\tA$ if there are $X_3, Z_0, Z_1$ satisfying the following conditions such that the $\tB$ equals 
\[\left(\begin{bmatrix} X_0 \\ X_1 \\ X_2 \\ X_3 \end{bmatrix}, \begin{bmatrix} X_0 \\ X_1 \\ X_2 \\ X_3 \end{bmatrix}, \begin{bmatrix} Z_0 \\ Z_1\end{bmatrix}\right) \circ \tA.\]



\begin{enumerate}
\item $X_3$ is a matrix such that the row span of 
\[\begin{bmatrix} X_0 \\ X_1 \\ X_2 \\ X_3 \end{bmatrix}\]
is $(\Z / p^e \Z)^n$, and the rows of $X_3$ is a smallest generating set of 
\[\left\langle \left\{x \in O_{\mathfrak{X}, p^e} : \forall s \in S, s \cdot x^t = 0 \text{ and }  x \cdot \mathfrak{Z}_{\tA, z} \in S 
\text{ for any } z \in (\Z / p^e \Z)^m \text{ s.t. } X_1  \mathfrak{Z}_{\tA, z} X_1^t = 0 \right\}\right\rangle,\]
where \[S = \langle \{s  \in (\Z / p^e \Z)^m : \exists z  \text{ s.t. }X_1 \mathfrak{Z}_{\tA, z} X_1^t = 0, s \in \rowspan(T \mathfrak{Z}_{\tA, z})\}\rangle.\]

%\[\langle \{x \in O_{T, \mathfrak{X}, p^e} : \forall s \in \langle S\rangle, s \cdot x^t = 0 \text{ and } 
%x \cdot \mathfrak{Z}_{T, z} \in \langle S\rangle \text{ for all } \mathfrak{Z}_{T, z} \text{ s.t. } \begin{bmatrix} U_0 \\ U_1 \end{bmatrix}  \mathfrak{Z}_{T, z} \begin{bmatrix} U_0 \\ U_1 \end{bmatrix} ^t = 0 \}\rangle.\]
\item Suppose the matrices are ordered by the lexical order. Compute $z_1, \dots, z_d$ as row vectors such that \[\begin{bmatrix} X_0 \\ X_1 \\ X_2\end{bmatrix} \mathfrak{Z}_{T, z_1} \begin{bmatrix} X_0 \\ X_1 \\ X_2 \end{bmatrix}^t, \dots, \begin{bmatrix} X_0 \\ X_1 \\ X_2 \end{bmatrix} \mathfrak{Z}_{T, z_d} \begin{bmatrix} X_0 \\ X_1 \\ X_2 \end{bmatrix}^t\] is the smallest generator of \[\left\langle \left\{\begin{bmatrix} X_0 \\ X_1 \\ X_2 \end{bmatrix} \mathfrak{Z}_{T, z} \begin{bmatrix} X_0 \\ X_1  \\ X_2 \end{bmatrix}^t: z \in (\Z / p^{e_0} \Z)^m\right\}\right\rangle.\]
$Z_0$ is the matrix with $z_i$ as $i$-th row.
\item $Z_1$ is a matrix such that the rows of $Z_1$ is a smallest generating set of \[\left\langle \left\{z : \begin{bmatrix} X_0 \\ X_1  \\ X_2 \end{bmatrix} \mathfrak{Z}_{T, z_i} \begin{bmatrix} X_0 \\ X_1 \\ X_2 \end{bmatrix}^t = 0\right\} \right\rangle. \]
%\item Let $U = \begin{bmatrix} U_0 \\ U_1 \\ U_2 \end{bmatrix}$ and $Z = \begin{bmatrix} Z_0 \\ Z_1 \end{bmatrix}$. 
\end{enumerate}
%We say a semi-canonical form is with respect to parameters $n, m, u, v$, where $u$ is the number of rows for $U_0$ plus the number of rows for $U_1$, and $v$ is the number of rows for $Z_0$. 
\end{definition}

\begin{lemma}\label{lem:tensor_isomorphism_restricted}
Let $\tA$ and $\tB$ be two $n \times n \times m$ alternating tensors over $\Z / p^e \Z$. 
Suppose $\tA'$ is a semi-canonical form for $\tA$ obtained for $T_{\tA}, X_{\tA, 0}, X_{\tA, 1}, X_{\tA, 2}$,
 and $\tB'$ is a semi-canonical form for $\tB$ obtained for $T_{\tB}, X_{\tB, 0}, X_{\tB, 1}, X_{\tB, 2}$.
If $T_{\tA}, X_{\tA, 0}, X_{\tA, 1}, X_{\tA, 2}$ has the same number of rows to $T_{\tB}, X_{\tB, 0}, X_{\tB, 1}, X_{\tB, 2}$ respectively, then 
\begin{enumerate}
\item If there are invertible matrices $X$ and $Y$ such that 
the following conditions hold:
\begin{enumerate}
\item  $\tB = (X, X, Y) \circ \tA$;
\item $X_{\tB, 0} = X_{\tA, 0} X^{-1} + W_{0}$ for some $W_0$ of the same dimension as $X_{\tA, 0}$;
\item $X_{\tB, 1} = X_{\tA, 1} X^{-1} + W_{1}$ for some $W_1$ of the same dimension as $X_{\tA, 1}$; 
\item $X_{\tB, 2} = X_{\tA, 2} X^{-1} + W_{2}$ for some $W_2$ of the same dimension as $X_{\tA, 2}$; 
\item $T_{\tB} = T_{\tA} X^{-1} + W_{3}$ for some $W_3$ of the same dimension as $T_{\tA}$;
\item  For each $i \in \{0, 1, 2, 3\}$,  $(W_i, I, I) \circ \tB$ is a zero tensor,
\end{enumerate}
then there are invertible matrices
\[P = \begin{bmatrix} I_{n'} & 0 \\ 0 & P_0\end{bmatrix}\]
for some invertible matrix $P_0$
and 
\[Q = \begin{bmatrix} I_{m'} & Q_0 \\ 0 & Q_1\end{bmatrix}\]
for some arbitrary matrix $Q_0$ and invertible matrix $Q_1$
such that $\tB' = (P, P, Q)\circ \tA'$, 
where $n'$ is the sum of the numbers of rows for $X_{\tA, 0}, X_{\tA, 1}$, and $X_{\tA, 2}$,
and $m'$ is the number of rows of $Z_{\tA, 0}$ when constructing the semi-canonical form $\tA'$. 



%there are $T_{\tA}, X_{\tA, 0}, X_{\tA, 1}, X_{\tA, 2}$ for $\tA$ and $T_{\tB}, X_{\tB, 0}, X_{\tB, 1}, X_{\tB, 2}$ for $\tB$ such that 
%letting $\tA'$ be a semi-canonical form for $\tA$ obtained for $T_{\tA}, X_{\tA, 0}, X_{\tA, 1}, X_{\tA, 2}$
% and $\tB'$ be a semi-canonical form for $\tB$ obtained for $T_{\tB}, X_{\tB, 0}, X_{\tB, 1}, X_{\tB, 2}$,
%there are invertible matrices $P$ and $Q$ such that 

%and $(T_{\tA}, I, I)\circ \tA = (T_{\tB} X, X, Y) \circ \tB$
\item If there do not exist invertible matrices $P \in \GL(n, \Z / p^e \Z)$ and $Q \in \GL(m, \Z / p^e \Z)$ such that $\tB = (P, P, Q) \circ \tA$, 
then 
there are no invertible matrices $P'$ and $Q'$ such that $\tB' = (P', P', Q') \circ \tA'$.
\end{enumerate}

\end{lemma}
\begin{proof}
Let $X_{\tA, 3}, Z_{\tA, 0}$, and $Z_{\tA, 1}$ denote the matrices used to construct $\tA$ for $\tA$, and $X_{\tB, 3}, Z_{\tB, 0}, and Z_{\tB, 1}$ denote the matrices used to construct $\tA$ for $\tA$. 

For the first property, by Definition~\ref{def:semi_canonical}, 
we have $X_{\tB, 3} = X_{\tA, 3} X^{-1}$, $Z_{\tB, 1} = Z_{\tA, 1} Y^{-1}$, and 
$Z_{\tB, 0} = Z_{\tA, 0} Y^{-1} + U$ for some $U$ such that the row span of $U$ is a subset of the row span of $Z_{\tB, 1}$. 
Hence, matrices $P$ and $Q$ exist. 

The second property is obtained by the first condition of Definition~\ref{def:semi_canonical} and the fact that 
\[\rowspan\left(\begin{bmatrix} Z_{\tA, 0} \\ Z_{\tA, 1} \end{bmatrix}\right)  = (\Z / p^e \Z)^m \text{ and } \rowspan\left(\begin{bmatrix} Z_{\tB, 0} \\ Z_{\tB, 1} \end{bmatrix}\right)  = (\Z / p^e \Z)^m \] 
by the last condition of Lemma~\ref{lem:enumeration} and the definition of $Z_{\tA, 0}, Z_{\tA, 1}, Z_{\tB, 0}$, and $Z_{\tB, 1}$.  
\end{proof}

\begin{definition}
In the correlated matrix tuple isometry problem, the inputs contains 
two $n \times n$ alternating matrix tuples $(A_1, \dots, A_k)$ and $(A_1', \dots, A_k')$ over $ \Z / p^e \Z$,
two $m \times n$ matrix tuples $(B_1, \dots, B_\ell)$ and $(B_1', \dots, B_\ell')$ over $\Z / p^e \Z$,
and a parameter $t$. 
The problem is to determine if there is an invertible matrix $P \in \GL(n, \Z / p^e \Z)$
and another matrix 
\[Q = \begin{bmatrix} I_{t} & Q_0 \\ 0 & Q_1\end{bmatrix}\] for some $Q_0 \in M(t, m - t, \Z / p^e \Z)$, $Q_1 \in M(m - t, \Z / p^e \Z)$ such that 
$A_i' = P A_i P^t$ for every $1 \leq i \leq k$ and $B_i' = Q B_i P^t$ for every $1 \leq i \leq \ell$.
\end{definition}
\begin{lemma}\label{lem:correlated_matrix_tuple}
The correlated matrix tuple isometry problem reduces to the alternating matrix tuple isometry problem in polynomial time in the following sense: 
For an input instance of the correlated matrix tuple isometry problem containing 
$n \times n$ alternating matrix tuples $(A_1, \dots, A_k)$ and $(A_1', \dots, A_k')$ over $ \Z / p^e \Z$ and 
$m \times n$ matrix tuples $(B_1, \dots, B_\ell)$ and $(B_1', \dots, B_\ell')$ over $\Z / p^e \Z$,
there is an algorithm to construct alternating $g \times g$ matrix tuples 
$(C_1, \dots, C_h)$ and $(C_1, \dots, C_h)$ over $\Z / p^e \Z$ with $g = \mathrm{poly}(n, m)$ and $h =  \mathrm{poly}(k, \ell)$
such that $(C_1, \dots, C_h)$ and $(C_1, \dots, C_h)$ are isometric if and only if 
the input of the correlated matrix tuple isometry  problem 
$(A_1, \dots, A_k)$  and $(B_1, \dots, B_\ell)$ are correlated isometric to $(A_1', \dots, A_k')$ and $(B_1', \dots, B_\ell')$ respectively. 
%\ynote{The statement is missing here.} \xnote{Fixed}
\end{lemma}

A proof of Lemma~\ref{lem:correlated_matrix_tuple}  for matrix tuples over a finite field is provided in Section III.D of~\cite{IMQSZ24}. The proof for matrix tuples over the ring 
$\mathbb{Z}/p^e\mathbb{Z}$ is identical to the finite field case, so we omit it here.


Finally, we give our algorithm in Figure~\ref{fig:algorithm} and prove Theorem~\ref{thm:tensor_iso_algo}. 

\begin{figure}[th]

\begin{framed}
\noindent \textbf{Alternating Matrix Space Isometry Algorithm}

\noindent \textbf{Input}: Two tensors $\tA$ and $\tA'$ of the same dimension $n \times n \times m$ over $\Z / p^e \Z$.

\noindent \textbf{Output}: Yes or no.

\begin{enumerate}
\item If the parameters $\alpha_1, \dots, \alpha_e, \beta_1, \dots, \beta_e, \gamma_1, \dots, \gamma_e$ of the two tensors are not identical, then return no. 
Otherwise, set  $e_0 = \lceil (\sum_{\ell = 1}^{e} \ell \cdot \beta_\ell)^{1/9} \rceil$, and $r = \lceil ((\sum_{\ell = 1}^{e} \ell \cdot \alpha_\ell) \cdot \log p)^{2/9} \rceil$. 
%\item Let $r = ???$ and $e = ???$. 
\item Find arbitrary matrices $X_{\tA, 0}$, $X_{\tA, 1}, X_{\tA, 2}, T_\tA$ satisfying the conditions of Lemma~\ref{lem:enumeration}, and compute the semi-canonical form $\tA'$ for $\tA$ with respect to $X_{\tA, 0}$, $X_{\tA, 1}, X_{\tA, 2}, T_\tA$ using arbitrary matrices $X_{\tA, 3}, Z_{\tA, 0}, Z_{\tA, 1}$ satisfying the conditions of Lemma~\ref{lem:enumeration}. 
\item 
Let $n'$ and $m'$ be the integers such that $\tA'$ is of dimension $n' \times n' \times m'$.  
Let $n''$ be the sum of the numbers of rows for $X_{\tA, 0}$, $X_{\tA, 1}$, and $X_{\tA, 2}$, and $m''$ be the number of rows for $Z_{\tA, 0}$. 
\item For any $X_{\tB, 0}$, $X_{\tB, 1}, X_{\tB, 2}, T_\tB$ for $\tA'$ satisfying the conditions of Lemma~\ref{lem:enumeration} with dimensions the same as $X_{\tA, 0}$, $X_{\tA, 1}, X_{\tA, 2}, T_\tA$ respectively,  
\begin{enumerate}
\item Construct a semi-canonical form $\tB'$ for $\tB$ with respect to $X_{\tB, 0}$, $X_{\tB, 1}, X_{\tB, 2}, T_\tB$. 
%\item Let $U_2$ be the matrix such that the rows of $U_2$ is a smallest generating set of  $\langle \{x \in O_{T, \mathfrak{X}, p^e} : \forall s \in \langle S\rangle, s \cdot x^t = 0 \text{ and } x \cdot \mathfrak{Z}_{T, z} \in \langle S\rangle \text{ for all } \mathfrak{Z}_{T, z} \text{ s.t. } \begin{bmatrix} U_0 \\ U_1 \end{bmatrix}  \mathfrak{Z}_{T, z} \begin{bmatrix} U_0 \\ U_1 \end{bmatrix} ^t = 0 \}\rangle$, and  $U_2'$ be the matrix such that the rows of $U_2'$ is a smallest generating set of  $\langle \{x \in O_{T', \mathfrak{X}, p^e} : \forall s' \in \langle S'\rangle, s' \cdot x^t = 0 \text{ and } x \cdot \mathfrak{Z}_{T', z} \in \langle S'\rangle \text{ for all }\mathfrak{Z}_{T', z} \text{ s.t. } \begin{bmatrix} U_0' \\ U_1' \end{bmatrix}  \mathfrak{Z}_{T', z} \begin{bmatrix} U_0' \\ U_1' \end{bmatrix} ^t = 0 \}\rangle$.
%\item Compute $z_1, \dots, z_d$ as row vectors such that $\begin{bmatrix} U_0 \\ U_1 \end{bmatrix} \mathfrak{Z}_{T, z_1} \begin{bmatrix} U_0 \\ U_1 \end{bmatrix}^t, \dots, \begin{bmatrix} U_0 \\ U_1 \end{bmatrix} \mathfrak{Z}_{T, z_d} \begin{bmatrix} U_0 \\ U_1 \end{bmatrix}^t$ is the smallest generator of $\langle \{\begin{bmatrix} U_0 \\ U_1 \end{bmatrix} \mathfrak{Z}_{T, z} \begin{bmatrix} U_0 \\ U_1 \end{bmatrix}^t: z \in (\Z / p^{e_0} \Z)^m\}\rangle$ suppose all the matrices are ordered by their lexical order.   Compute $z_1', \dots, z_{d'}$ as row vectors such that $\begin{bmatrix} U_0' \\ U_1' \end{bmatrix} \mathfrak{Z}_{T', z_1'} \begin{bmatrix} U_0' \\ U_1' \end{bmatrix}^t, \dots, \begin{bmatrix} U_0' \\ U_1' \end{bmatrix} \mathfrak{Z}_{T', z_{d'}'} \begin{bmatrix} U_0' \\ U_1' \end{bmatrix}^t$ is the smallest generator of $\langle \{\begin{bmatrix} U_0' \\ U_1' \end{bmatrix} \mathfrak{Z}_{T', z} \begin{bmatrix} U_0' \\ U_1' \end{bmatrix}^t: z \in (\Z / p^{e_0} \Z)^m\}\rangle$ suppose all the matrices are ordered by their lexical order.  
%\item If $U_2$ and $U_2'$ have different number of rows, or $d \neq d'$, then continue. 
%\item If $\begin{bmatrix} U_0 \\ U_1 \end{bmatrix} \mathfrak{Z}_{T, z_i} \begin{bmatrix} U_0 \\ U_1 \end{bmatrix}^t \neq \begin{bmatrix} U_0' \\ U_1' \end{bmatrix} \mathfrak{Z}_{T', z_i'} \begin{bmatrix} U_0' \\ U_1' \end{bmatrix}^t$ for some $i$, then continue.
\item If $\tA'_{i, j, k} \neq \tB'_{i, j, k}$ for any $1 \leq i \leq n'', 1 \leq j \leq n'', 1 \leq k \leq m''$, then continue. 
%\item $Z_0$ be the matrix with $z_i$ as $i$-th row, and $Z_1$ be the matrix with rows as a smallest generating set of $\langle z : \begin{bmatrix} U_0 \\ U_1 \end{bmatrix} \mathfrak{Z}_{T, z_i} \begin{bmatrix} U_0 \\ U_1 \end{bmatrix}^t = 0\rangle$. Let $Z_0', Z_1'$ be matrices constructed in the same way for $T'$. 
\item Let $A_i$ be the induced matrix of $\mathfrak{Z}_{\tA, i}$ on the last $n' - n''$ rows and columns for any $1 \leq i \leq m''$, and $B_i$ be $\mathfrak{X}_{C, i}$ removing the first $n''$ columns for any $1 \leq i \leq n''$. Construct $A_i'$ and $B_i'$ for $\tB'$ in the same way. 
\item If $A_1, \dots, A_{m''}, B_1, \dots, B_{n''}$ and $A_1', \dots, A_{m''}', B_1', \dots, B_{n''}'$ are correlated matrix tuple isometry, then output yes. 
\end{enumerate}
\item Output no.
\end{enumerate}
\end{framed}
\caption{Alternating Matrix Space Isometry Algorithm}\label{fig:algorithm}
\end{figure}
%\begin{lemma}\label{lem:tensor_iso_algo_main}
%There is an algorithm to determine whether two $n \times n \times m$ tensors over $\Z / p^{e_0} \Z$ in time $p^{(\sum_{\ell = 1}^{e} \ell \cdot (\beta_\ell + \alpha_\ell)) \cdot(\sum_{\ell = 1}^{e} \ell \cdot (\beta_\ell + \alpha_\ell)\log p)^{8/9}}$. 
%\end{lemma}

\begin{proof}[Proof of Theorem~\ref{thm:tensor_iso_algo}]
The correctness of the algorithm is obtained by Lemma~\ref{lem:enumeration}, Lemma~\ref{lem:lexical_order}, Definition~\ref{def:semi_canonical},  Lemma~\ref{lem:tensor_isomorphism_restricted}, and Lemma~\ref{lem:correlated_matrix_tuple}.

In the rest of this proof, we bound the running time of the algorithm. When enumerating rows of $X_{\tA, 0}$ and $X_{\tB, 0}$, we only enumerate row vectors $x$ such that $\mathfrak{X}_{T, x}$ are different. 
By Lemma~\ref{lem:enumeration}, the cost of enumerating a single row of $X_0$ is $p^{\sum_{\ell = 1}^{e} \ell \cdot \beta_\ell}$, and the cost of enumerating $X_0$ is 
\[\left(p^{(\sum_{\ell = 1}^{e} \ell \cdot \beta_\ell)} \right)^{(\sum_{\ell = e_0 + 1}^{e} \beta_\ell)}.\]


Similarly, when enumerating 
$X_{\tA, 1}, X_{\tA, 2}, T_{\tA}, X_{\tB, 1}, X_{\tB, 2}$, or $T_{\tB}$, we only enumerate row vectors $x$ such that $\mathfrak{X}_{T, x}$ are different.
By Lemma~\ref{lem:enumeration},
the cost of enumerating a single row of $X_{\tA, 1}, X_{\tA, 2}, T_{\tA}$, $X_{\tB, 1}$, $X_{\tB, 2}$, or $T_{\tB}$ is \[p^{\sum_{\ell = 1}^{e_0} \ell \cdot \beta_\ell} \cdot p^{\sum_{\ell = e_0 + 1}^{e} e_0 \cdot \beta_\ell},\] and thus the cost of enumerating all possible $X_{\tA, 1}, X_{\tA, 2}, T_{\tA}, X_{\tB, 1}, X_{\tB, 2}$, and $T_{\tB}$  is 
\[\left(p^{(\sum_{\ell = 1}^{e_0} \ell \cdot \beta_\ell + \sum_{\ell = e_0 + 1}^{e} e_0 \cdot \beta_\ell)}\right)^{ O(\max\{\log(p^{\sum_{i = 1}^{e} i \cdot \alpha_i}), r\}/ \sqrt{r} + e_0^2 r^2 + \sum_{i = e_0 + 1}^{e} \beta_i)}.\]

%The enumeration of $\langle S \rangle$ is equivalent to enumerate at most $O(e^2 r^2 + \sum_{i = e + 1}^{e_0} \beta_i)$ row vectors in $O_{\mathfrak{X}, p^e}$, whose cost is 
%\[p^{(\sum_{\ell = 1}^{e} \ell \cdot \beta_\ell + \sum_{\ell = e + 1}^{e_0} e \cdot \beta_\ell)\cdot O(e^2 r^2 + \sum_{i = e + 1}^{e_0} \beta_i)}.\]

By the choices of the parameters, 
%Let $e = \lceil (\sum_{\ell = 1}^{e_0} \ell \cdot \beta_\ell)^{1/9} \rceil$, and $r = \lceil ((\sum_{\ell = 1}^{e_0} \ell \cdot \alpha_\ell) \cdot \log p)^{2/9} \rceil$. 
the total cost of enumeration is at most 
\[p^{O((\sum_{\ell = 1}^{e} \ell \cdot (\beta_\ell + \alpha_\ell)) \cdot(\sum_{\ell = 1}^{e} \ell \cdot (\beta_\ell + \alpha_\ell)\log p)^{8/9})}.\]
Recall that $\mu = \sum_{\ell = 1}^e (\alpha_\ell + \beta_\ell + \gamma_\ell)$, %\ynote{Added:} 
where $\gamma_i=\beta_i$ here as we work with alternating tensors. 
So the total cost of enumeration is at most 
\[(p^\mu)^{O(( \mu \cdot \log p)^{8/9})}. \qedhere\]
\end{proof}
\begin{comment}
\begin{lemma}
There exists an algorithm that determines the isomorphism of two nilpotent class-2 groups of odd order in time $N^{(\log N)^{8/9}}$, where $N$ is the order of the two groups.
\end{lemma}
\begin{proof}
Construct the tensors for the two groups. Then we know $G / Z(G)$ is isomorphic to 
\[\bigoplus_{\ell = 1}^{e_0} \left(\Z / p^{\ell } \Z\right)^{\beta_\ell},\]
and $[G, G]$ is isomorphic to 
\[\bigoplus_{\ell = 1}^{e_0} \left(\Z / p^{\ell } \Z\right)^{\alpha_\ell}.\]
Hence, we have $N \geq |G / Z(G)| \cdot |[G, G]| \geq p^{\sum_{\ell = 1}^{e_0} \ell \alpha_\ell} \cdot p^{\sum_{\ell = 1}^{e_0} \ell \beta_\ell} $. By Lemma~\ref{lem:tensor_iso_algo_main}, we prove this lemma. 
\end{proof}
\end{comment}


% !TEX root = TI5.tex

\subsection{Efficient algorithm for isometry of tuples of (anti)symmetric bilinear forms over $\Z/p^e \Z$ 
	%and finite commutative local PIRs of odd characteristic
}\label{subsec:tuple-isometry}\label{sec:matrix_tuple}

Throughout Section~\ref{subsec:tuple-isometry}, $R$ denotes a finite local commutative principal ideal ring (PIR) of odd characteristic (equivalently, $R/\m$ is a field of odd characteristic), and we use the notation and terminology laid out in Section~\ref{sec:PIR}.

Given an $R$-module $M$, integers $1 \leq z_i \leq e$, and $R$-bilinear forms $B_i,C_i \colon M \otimes_R M \to R/\m^{z_i}$, the \algprobm{(Tuple) Isometry Problem} is to decide whether there exists $F \in \aut(M)$ such that $B_i(F(u), F(v)) = C_i(u,v)$ for all $u,v \in M$ and all $i=1,\dotsc,m$, and any such $F$ is called an \emph{isometry} from $\vec{B}$ to $\vec{C}$. If we represent each bilinear form by its corresponding matrix, which we denote by the same letter by abuse of notation, we have $(B_i)_{jk} = B_i(v_j, v_k)$, where the $v_*$ are a basis of $M$, and the congruence problem then asks whether $F^t B_i F = C_i$ for all $i = 1,\dotsc,m$, where now we view $F$ as its matrix according to Theorem~\ref{thm:HillarRheaPIR}. (Note that the matrix $B_i$ can be any matrix with entries in $R/\m^{z_i}$, subject to the constraint that if $z_i \leq \min\{e_j, e_k\}$, then the $(j,k)$ entry must be a multiple of $\m^{z_i - \min\{e_j, e_k\}}$. This follows from bilinearity and the fact that $v_j$ is annihilated by $\m^{e_j}$ and $v_k$ is annihilated by $\m^{e_k}$ in $M$.)

For $\varepsilon \in \{\pm 1\}$, we say a bilinear form $B \colon M \otimes M \to R/\m^z$ is $\varepsilon$-symmetric if $B(u,v) = \varepsilon B(v,u)$ for all $u, v \in P$; thus $\varepsilon=1$ corresponds to symmetric and $\varepsilon=-1$ corresponds to anti-symmetric (equivalently, since $p$ is odd, alternating). 

Our main result in this section is:

\begin{theorem}[{Generalizes \cite[Thm.~1.7(1)]{IQ19} from finite fields to finite local commutative PIRs}]\label{thm:matrix_tuple_isometry}
If $R$ is a finite local commutative PIR with quotient field of odd characteristic $p$, and $M$ is an $n$-generated $R$-module, the \algprobm{Tuple Isometry Problem} for $m$-tuples of $\varepsilon$-symmetric bilinear forms on $M$ can be solved in randomized  time $\poly(n,m,e)$ with operations in $R$ at unit cost, or deterministic time $\poly(n,m,e, p)$ with operations in $R$ at unit cost, where $p$ is the characteristic of $R/\m$.

In particular, when $R = \Z/p^e \Z$, the latter implies a solution in deterministic time $\poly(n,m,p,e)$ on a Turing machine.
\end{theorem}

Our algorithm follows that of Ivanyos and Qiao closely, and uses several of their key results as subroutines. Here we recall the main steps of their algorithm, and give proofs of the key results that need to be generalized from the case of finite fields to local commutative PIRs as above; we maintain the numbering from their paper for ease of comparison. In Appendix~\ref{app:similar}, we include those proofs of the key results here that are identical or nearly identical to the original ones, merely so that it may be checked carefully that the stated proofs do indeed go through in our more general setting.

Throughout the remainder of this section, let $M$ be an $n$-generated $R$-module. Given an $\varepsilon$-symmetric bilinear form $B \colon M \otimes M \to R/\m^z$, we define its kernel as $\ker B = \{m \in M | B(m,m') = 0 \forall m' \in M\}$. (Note that since $B$ is $\varepsilon$-symmetric, its left kernel and right kernel coincide as submodules of $M$.) Given $z_1, \dotsc, z_m$ and $\varepsilon$-symmetric bilinear forms $B_i \colon M \otimes M \to R/\m^{z_i}$, we say the tuple $(B_1, \dotsc, B_m)$ is \emph{non-degenerate} if $\bigcap_{i=1}^m \ker B_i = 0$. 

%[Move this note to where we actually need it. May need to put something in the appendix for commutative local PIRs beyond Abelian groups... We note that linear algebra over a local commutative PIR can be done in polynomial time with operations in $R$ at unit cost, nearly as efficiently as over fields, in particular this holds (even on usual Turing machines) when $R = \Z/p^e \Z$ \cite{GoldmannRussell}.]

\paragraph{Steps of the algorithm.}
\begin{enumerate}
\item \textit{Reduce to the non-degenerate case.} 

\begin{lemma}
There is a polynomial-time $R$-algorithm which, given a tuple of $\varepsilon$-symmetric bilinear forms $\vec{B}$ on $M$, produces a quotient module $M'$ of $M$, and a non-degenerate tuple of $\varepsilon$-symmetric bilinear forms $\text{nondeg}(\vec{B})$ on $M'$ such that for any two tuples $\vec{B}, \vec{C}$, we have that $\text{nondeg}(\vec{B})$ and $\text{nondeg}(\vec{C})$ are isometric if and only if $\vec{B}, \vec{C}$ are isometric.
\end{lemma}

Over fields, this is easy by standard linear algebra. Over finite local PIRs, we instead use Smith Normal Form.

Essentially, $\vec{B}$ and $\text{nondeg}(\vec{B})$ are isometric---there is no real loss of information in moving from one to the other---but because they are bilinear forms on different modules $M$ and $M'$, respectively, making this precise takes some work. As we won't need such a statement, we avoid this complication, but we mention it as it may still be helpful for intuition.

\begin{proof}
We show how to compute $\text{nondeg}(\vec{B})$ and the associated quotient $M'$ of $M$; the same procedure is then applied to $\vec{C}$. (If the quotient module produced in computing $\text{nondeg}(\vec{C})$ is not the same as $M'$, then $\vec{B}$ and $\vec{C}$ do not have isomorphic kernels, and therefore the two tuples cannot be isometric.)

Suppose $\vec{B}$ is degenerate, with common kernel $\bigcap_{i=1}^m \ker B_i = K \leq M$. By Smith Normal Form, there is a basis $m_1,\dotsc,m_n$ for $M$, together with exponents $0 \leq k_1 \leq k_2 \leq \dotsb $ and an integer $k$ such that $t^{k_1} m_1, t^{k_2} m_2, \dotsc, t^{k_k} m_k$ is a basis for $K$. By our assumptions on $R,M,\vec{B}$, we can compute this SNF in polynomial time, which includes the elements $m_i$ and the exponents $k_i$. If the $B_i$ were originally written in a different basis, we apply the change of basis matrix (coming from the SNF computation) to put $B_i$ into the $\{m_1, \dotsc, m_n\}$ basis, resulting in an isometric tuple of bilinear forms, which we call $\vec{B}'$.

In this basis, for $j=1,\dotsc,k$, we have that $t^{k_j} R / \m^{e_j}$ (in the $j$-th position) is in the kernel of $\vec{B}'$. Since every such vector is in the kernel of each $B_i'$, it must be the case that every element of the $j$-th row and $j$-th column of $B_i'$ is divisible by $t^{e_j - k_j}$. 


Let $M' = M / K$. If $M = \bigoplus_{i=1}^n \langle m_i \rangle = \bigoplus_{i=1}^n R/\langle t^{e_i} \rangle$, then we have $M'$ is the following ``coordinate-wise quotient'': $M' = \left(\bigoplus_{i=1}^k R/\langle t^{k_i}  \rangle\right) \oplus \left(\bigoplus_{i=k+1}^n R/\langle t^{e_i} \rangle \right)$. That is, in the first $k$ coordinates, we mod out by $t^{k_i}$ in the $i$-th coordinate, and the remaining coordinates are untouched. (If some $k_i=0$, we may then omit that coordinate, which corresponds to $M'$ having a smaller basis than $M$.)

In $\text{nondeg}(B)$ we replace each element of the $j$-th row and column of $B'$ with its quotient by $t^{e_j - k_j}$. (Although there are multiple such quotients, they all differ by multiples of $t^{k_j}$, hence they are all equal as elements of $R/\langle t^{k_j} \rangle$, which is how we now view the $j$-th row and column of $\text{nondeg}(B)$.) If we ever have $e_j = k_j$, then that row and column of every $B_i'$ is zero, and we simply remove that row and column from each $B_i'$. For $j=1,\dotsc,k$, since the $j$-th row and column of $\text{nondeg}(B)$ are now viewed as elements of $R/t^{k_j}$, the matrices we constructed in $\text{nondeg}(\vec{B})$ correspond to bilinear forms on $M'$, whose common kernel corresponds to the image of $K$ in $M'$, which is zero. Thus $\text{nondeg}(\vec{B})$ is non-degenerate.

Furthermore, if $\vec{B}$ and $\vec{C}$ are isometric, then so are $\vec{B}'$ and $\vec{C}'$ (since $\vec{B}'$ is isometric to $\vec{B}$ by construction, and similarly for $\vec{C}'$ and $\vec{C}$). If $F$ is an isometry from $\vec{B}'$ to $\vec{C}'$, let $F'$ be the matrix gotten from $F$ by removing the rows and columns corresponding to rows and columns of $\vec{B}'$ that are removed in $\text{nondeg}(\vec{B})$. Then $F'$ is an isometry from $\text{nondeg}(\vec{B})$ to $\text{nondeg}(\vec{C})$. Doing this process in reverse gives the converse.
\end{proof}

From now on, we assume that $\vec{B}$ and $\vec{C}$ are themselves non-degenerate tuples of bilinear forms on $M$.

\item \textit{Solve the twisted equivalence problem.} Test whether $\vec{B}$ and $\vec{C}$ are ``twisted equivalent'', meaning that there are invertible elements $A,D \in \aut(M)$ such that $A^t \vec{B} = \vec{C} D$, or in terms of bilinear forms, that $B_i(Au, v) = C_i(u, Dv)$ for all $u,v \in M$ and all $i=1,\dotsc,m$.

\begin{lemma} \label{lem:twisted}
There is a polynomial-time reduction from \algprobm{Twisted Equivalence} of tuples of $\varepsilon$-symmetric bilinear forms on an $R$-module $M$ to the \algprobm{Module Isomorphism Problem} over $R$. In particular, \algprobm{Twisted Equivalence} for such tuples can be solved in polynomial time.
\end{lemma}

The reduction is the same construction as in Ivanyos \& Qiao \cite[Prop.~3.2]{IQ19}, but the argument of correctness requires some care in the case of a finite local PIR, so we include it in Appendix~\ref{app:similar}. To solve the \algprobm{Module Isomorphism Problem} over $R$, we use the polynomial-time algorithm of Ciocanea-Teodorescu \cite{CT}, which extends that of Brooksbank \& Luks \cite{BL08} from working over fields to working over finite rings.

\item \textit{Reduce to decomposing an element of a $*$-algebra.}

\newcommand{\adj}{\text{Adj}}

\begin{definition}[{cf. \cite[Def.~2.3]{IQ19}}]
Let $R$ be a finite local PIR with quotient field of odd characteristic, and $M$ a finite $R$-module. For $\varepsilon \in \{\pm 1\}$, and $\vec{B} = (B_1, \dotsc, B_m)$ a tuple of $\varepsilon$-symmetric bilinear forms on $M$, the \emph{adjoint algebra} of $\vec{B}$, denoted $\adj(\vec{B})$, is $\{(A,D) : A,D \in \End_R(M), A^t B_i = B_i D \,\forall i\}$. $\adj(\vec{B})$ is a $*$-algebra with $(A,D)^* = (D,A)$.
\end{definition}

By definition we have that $\adj(\vec{B})$ is a sub-algebra of $\End(M)^{op} \oplus \End(M)$, where $\End(M)^{op}$ has the same set of elements of $\End(M)$, but multiplication is reversed (corresponding to the use of the transpose in the above definition). However, when $\vec{B}$ is non-degenerate, we may realize $\adj(\vec{B})$ as isomorphic to a sub-algebra of $\End(M)$ as follows. The idea is that when $\vec{B}$ is non-degenerate, if $(A,D) \in \adj(\vec{B})$, then $A$ uniquely determines $D$ and vice versa. For suppose $(A,D), (A',D)$ are both in $\adj(\vec{B})$. Then we have $A^t B_i = B_i D = (A')^t B_i$ for all $i$. But then we have $(A^t - (A')^t)B_i = 0$ for all $i$, and thus the columns of $A - A'$ are in the common kernel of the $B_i$. But if $\vec{B}$ is non-degenerate, this common kernel is zero by definition, and thus $A - A' = 0$, so $A = A'$. From this it follows that the projection map from $\adj(\vec{B}) \subseteq \End(M)^{op} \oplus \End(M)$ onto the second coordinate $\End(M)$ is injective. Equivalently, when $\vec{B}$ is non-degenerate, we may thus define $\adj(\vec{B})$ as the following sub-algebra of $\End(M)$: $\adj(\vec{B}) \cong \{D \in \End(M) : (\exists A \in \End(M))[A^t \vec{B} = \vec{B} D]\}$, and the $*$ operation is then defined by $D^*$ being the unique $A$ such that $A^t \vec{B} = \vec{B} D$. With this setup, the remainder of \cite[Sec.~3.4]{IQ19} goes through without change. This is because the proofs there are phrased entirely in terms of the $*$ and inverse operations and their basic properties, not relying on anything special to working over fields. In particular, we have the following lemma:

\begin{lemma}[{Cf. \cite[Prop.~3.4]{IQ19}}]
Let $\vec{B}, \vec{C}$ be $\varepsilon$-symmetric bilinear forms on the $R$-module $M$, and suppose that $A,D \in \aut(M)$ are such that $A^t \vec{B} = \vec{C} D$. Let $E = A^{-1} D^{-1}$ and $\mathfrak{A} = \adj(\vec{C})$, viewed as a subalgebra of $\End(M)$ as above. Then $\vec{B}$ and $\vec{C}$ are isometric if and only if there exists $X \in \mathfrak{A}$ such that $X^* X = E$.
\end{lemma}

Given $E$ and $\mathfrak{A}$, finding such an $X$ is the \algprobm{$*$-Symmetric Decomposition Problem}. Note that here we are given $\mathfrak{A}$ by a set of elements of $\End(M)$ (which we may view as matrices over $R$ by Theorem~\ref{thm:HillarRheaPIR}) which form a basis for $\mathfrak{A}$.  

\item \textit{Solve the $*$-symmetric decomposition problem.}

\begin{enumerate}
\item \textit{Compute the algebra structure of $\mathfrak{A}$.} 

\begin{theorem}[{cf. \cite[Thm.~3.6]{IQ19}, which is based on \cite{Ron90,EG00,Iva00}}] \label{thm:structure}
Given an $n$-generated  $R$-module $M$, and a basis of a subalgebra $\mathfrak{A} \leq \End(M)$, there is a Las Vegas $R$-algorithm running in time $\poly(n)$ that computes (1) a basis of $\rad(\mathfrak{A})$, (2) a surjective $R$-algebra homomorphism $\pi \colon \mathfrak{A} \to M(n_1, \F_{q_1}) \oplus \dotsb \oplus M(n_k, \F_{q_k})$ with kernel $\rad(\mathfrak{A})$, where $\F_{q_i}$ is a finite extension field of $R/\m$ that is specified by a basis of $\F_{q_i}$ over $R/\m$. There is also a deterministic $R$-algorithm for the same running in time $\poly(n, p)$, where $p$ is the characteristic of $R/\m$.

Given the output of the aforementioned algorithm, there are also efficient deterministic algorithms to (3) compute $\pi(a)$ given $a \in \mathfrak{A}$, and (4) given $b \in M(n_1, \F_{q_1}) \oplus \dotsb \oplus M(n_k, \F_{q_k})$, compute $a \in \mathfrak{A}$ such that $\pi(a) = b$.
\end{theorem}

A key new ingredient compared to the field case is the following lemma:

\begin{lemma} \label{lem:rad}
Suppose $\mathfrak{A}$ is a finite (or Artinian) $R$-algebra, and let $\pi \colon \mathfrak{A} \to \mathfrak{A} / \m \mathfrak{A}$ be the natural quotient map. Then $\pi(\rad(\mathfrak{A})) = \rad(\mathfrak{A}/\m \mathfrak{A})$.
\end{lemma}

An element $x \in \mathfrak{A}$ is \emph{strongly nilpotent} if $xy$ is nilpotent for all $y$; in a finite (or Artinian) algebra, the Jacobson radical $\rad(\mathfrak{A})$ is equal to the set of all strongly nilpotent elements.

\begin{proof}
By Proposition 4.6 of Lam \cite{Lam}, in any ring $\mathfrak{A}$, if $I$ is an ideal contained in $\rad(\mathfrak{A})$, then $\rad(\mathfrak{A})/I = \rad(\mathfrak{A}/I)$. So it will suffice for us to show that $\m \mathfrak{A}$ is an ideal contained in $\rad(\mathfrak{A})$. If $x \in\m \mathfrak{A}$,%\ynote{Is it $\m pA$ or $\m A$?} 
then $x=tx'$ for some $x' \in A$, where $t$ is a generator of $\m$, and thus $xy=tx'y$. As $t$ is central, we have that $(xy)^e$ is divisible by $t^e$, hence zero. Thus every element of $\m \mathfrak{A}$ is strongly nilpotent, so $\m \mathfrak{A} \subseteq \rad(\mathfrak{A})$. Furthermore, since $t$ is central, $\m \mathfrak{A}$ is a (principal, two-sided) ideal. Thus $\m \mathfrak{A}$ is an ideal of $\mathfrak{A}$ contained in $\rad(\mathfrak{A})$, and the result follows.
\end{proof}

\begin{proof}[Proof of Theorem~\ref{thm:structure}] 
(1) To compute $\rad(A)$ we proceed as follows. If $\{a_1, \dotsc, a_d\}$ is a generating set of the additive group of $\mathfrak{A}$, then $\m \mathfrak{A}$ is spanned by the elements $\{ta_i : i = 1, \dotsc, d\}$. Furthermore, we can calculate a basis, exponents, and structure coefficients (as above) of $\mathfrak{A}/\m \mathfrak{A}$ by simply setting each $e_i=1$ and taking each $c_{ij}^k$ to its value modulo $t$. This also specifies for us the quotient map $\pi' \colon \mathfrak{A} \to \mathfrak{A}/\m \mathfrak{A}$ in terms of how it acts on a basis. Then using previous algorithms over finite fields \cite{Ron90,EG00,Iva00} in $\mathfrak{A}/\m \mathfrak{A}$---which is now a finite-dimensional algebra over the finite field $R/\m$---we can compute a basis for $\rad(\mathfrak{A}/\m \mathfrak{A})$ over $R/\m$, and for each basis element of $\rad(\mathfrak{A}/\m \mathfrak{A})$, take any one preimage under $\pi$. Those preimages, together with the aforementioned generating set of $\m \mathfrak{A}$, generate $\rad(\mathfrak{A})$ as an $R$-module, by Lemma~\ref{lem:rad}. Finally, using Smith Normal Form, we can extract from this generating set a basis of $\rad(\mathfrak{A})$.

(2 and 3) In the preceding step, we computed the quotient map $\pi' \colon \mathfrak{A} \to \mathfrak{A} / \m \mathfrak{A}$, the latter of which is an Artinian algebra over the field $R/\m$. Using previous algorithms over finite fields \cite{Ron90,EG00,Iva00}, we can then compute the quotient map $\mathfrak{A} / \m \mathfrak{A} \to \mathfrak{A} / \rad(\mathfrak{A})$ as well as the isomorphism of the latter with $M(n_1, \F_{q_1}) \oplus \dotsb \oplus M(n_k, \F_{q_k})$, where $\F_{q_i}$ is a finite extension field of $R/\m$ that is specified by a basis of $\F_{q_i}$ over $R/\m$.

(4) From the algorithms over finite fields \cite{Ron90,EG00,Iva00}, given $b \in M(n_1, \F_{q_1}) \oplus \dotsb \oplus M(n_k, \F_{q_k})$, we can compute an $a' \in \mathfrak{A} / \m \mathfrak{A}$ that is a preimage of $b$ under the quotient map $\mathfrak{A} \to \mathfrak{A} / \rad(\mathfrak{A})$. We now need to compute an $a \in \mathfrak{A}$ whose image modulo $\mathfrak{m}$ is $a'$. Let $M$ be the inclusion matrix of $\m \mathfrak{A}$ into $\mathfrak{A}$ (say, viewing ring elements here as column vectors). Putting $M$ into Smith Normal Form as in the first step, we get a basis $\{a_1,\dotsc,a_d\}$ of $\mathfrak{A}$ and exponents $f_1, \dotsc, f_d$ such that $t^{f_i} a_i = 0$ and $f_i$ is the least integer with this property, and furthermore we get a $k \leq d$ and integers $0 \leq e_1 \leq e_2 \leq \dotsb \leq e_k$, such that $\{t^{e_1} a_1, t^{e_2} a_2, \dotsb, t^{e_k} a_k\}$ is a basis for $\m \mathfrak{A}$. In this basis on $\mathfrak{A}$, the projection map $\mathfrak{A} \to \mathfrak{A} / \m \mathfrak{A}$ takes the form $\sum_{i=1}^d r_i a_i \mapsto \sum_{i=1}^k (r_i \text{ mod } t^{e_i}) \overline{a}_i$, where $\overline{a}_i$ represent the image of $a_i$ in $\mathfrak{A} / \m \mathfrak{A}$, and each $r_i$ is taken modulo the annihilator of $a_i$ (which is some power of $\m$). Now, say $a' = \sum_{i=1}^k \overline{r}_i \overline{a}_i$ is the element of $\mathfrak{A} / \m \mathfrak{A}$ whose preimage we want to compute. Then it suffices to lift $\overline{r}_i$ from $R/\m^{e_i}$ to $r_i \in R/\Ann_R(a_i)$, and then $\sum_{i=1}^k r_i a_i$ will be a preimage of $a'$.

By induction, it will suffice to be able to lift $\overline{r} \in R/\m^f$ to $r \in R/\m^{f+1}$ such that $r \equiv \overline{r} \pmod{\m^f}$, for arbtirary $0 \leq f \leq e$. However, this follows from our assumption that elements of $R/\m^f$ are represented by elements of $R$; see the beginning of Section~\ref{sec:PIR}.
\end{proof}

\item \textit{Recognize the $*$-algebra structure.} This part of the algorithm only deals with the semisimple quotient of $\mathfrak{A}$, which is necessarily an algebra over a field. Thus the results of \cite[Sec.~3.5.2]{IQ19} apply directly here, since this part of the algorithm is working in the same setting as they were working in, not a generalization thereof.

\item \textit{Reduce to the semi-simple case.} 

\begin{proposition}[{Cf. \cite[Prop.~3.11]{IQ19}}] \label{prop:semisimple}
Let $\mathfrak{A}$ be an $*$-algebra over $R$. Let $E \in \mathfrak{A}$ be an invertible $*$-symmetric element, and suppose there exists $Y \in \mathfrak{A} / \rad(\mathfrak{A})$ such that $Y^* Y \equiv E \pmod \rad(\mathfrak{A})$. Then there exists $X \in \mathfrak{A}$ such that $X^* X = E$, and such an $X$ can be found by a deterministic polynomial-time $R$-algorithm.
\end{proposition}

The one additional ingredient we need in the proof is the following lemma:

\begin{lemma}
Suppose $\mathfrak{A}$ is a finite $R$-algebra. Let $n = \dim_{R/\m} (\mathfrak{A} / \m \mathfrak{A})$. 
Then $(\rad \mathfrak{A})^{2^k} = 0$ for some $k \leq 1 + \lceil \log_2(\max\{n,e\}) \rceil $ (where recall that $e$ is minimal such that $\m^e = 0$).
\end{lemma}

\begin{proof}
By Lemma~\ref{lem:rad}, $\rad \mathfrak{A} \pmod \m = \rad(\mathfrak{A} / \m \mathfrak{A})$. As argued in the proof of \cite[Prop~3.11]{IQ19}, $\rad(\mathfrak{A}/\m\mathfrak{A})^{2^{k'}} = 0$ for some $k' \leq \lceil \log n \rceil$. Lifting to $\mathfrak{A}$, this is equivalent to saying that $\rad(\mathfrak{A})^{2^{k'}}$ is contained in $\m \mathfrak{A}$ for some $k' \leq \lceil \log n \rceil$. But since $\m^e = 0$ and $\m \subseteq R \subseteq \mathfrak{A}$ is central in $\mathfrak{A}$ (by the definition of $R$-algebra), we have that $(\m \mathfrak{A})^e = 0$ as well. Thus $\rad(\mathfrak{A})^{2^{k'} + e} = 0$. Finally, we have $2^{k'} + e \leq 2\cdot \max\{2^{k'}, e\} \leq 2^k$ for some $k \leq  1+ \lceil \log_2(\max\{n,e\}) \rceil$.
 \end{proof}

With this lemma in hand, the proof of Proposition~\ref{prop:semisimple} is nearly identical to that in Ivanyos--Qiao, again because their proof is phrased almost entirely in the language of $*$-algebras without using any special properties of algebras over fields. The bound on $k$ such that $(\rad \mathfrak{A})^{2^k} = 0$ is used in the proof to give an upper bound on the number of iterations of the algorithm, thus the running time of our $R$-algoirthm is still only at most $\poly(\dim_{R/\m} \mathfrak{A}/\m\mathfrak{A},e)$.

\item At this point, we may directly use the algorithm of \cite[Sec.~3.5.4 and 3.5.5]{IQ19} to solve the semi-simple case, since semi-simple algebras are always over fields. Thus the remainder of the algorithm uses their algorithm as a (black box) subroutine.
\end{enumerate}


\end{enumerate}


Finally, Theorem~\ref{thm:tensor_isomorphism_main} is obtained by Theorem~\ref{thm:GpI-to-AMSI}, Theorem~\ref{thm:tensor_iso_algo}, and Theorem~\ref{thm:matrix_tuple_isometry}. 

%\jnote{Notes for us: the places where we need $R$ to be finite---equivalently, $\m$ to be nilpotent---are in the reduction to the nondegenerate case above, and in Lemma~\ref{lem:rad}. Perhaps Lemma~\ref{lem:rad} can be proved without this assumption by other means, but I do not see how to reduce to the nondegenerate case when $R$ is not finite. e.g. in the $p$-adic integers, if the kernel includes the ideal $\langle p^i \rangle$, as a $\Z_p$-module, this is isomorphic to $\Z_p$...}


%[First recall how $A$ is given to us. We are given a "basis" $\{a_1,\dotsc,a_d\}$ (generating set of the additive group of $A$), exponents $e_i$ where $p^{e_i} a_i = 0$ and $e_i$ is minimal with this property, where the underlying additive group of $A$ is the direct sum of the cyclic groups generated additively by the $a_i$, say $A = \bigoplus_i \Z/p^{e_i} \Z$, and finally structure coefficients $c_{ij}^k$, where $a_i \cdot a_j = \sum_k c_{ij}^k a_k$, so $c_{ij}^k$ is in $\Z/p^{e_k} \Z$ for all $i,j,k$.]





